% =============================================================================
% E4 -- Workload-aware configuration selection (Tanish)
%
% Body only: % =============================================================================
% E4 -- Workload-aware configuration selection (Tanish)
%
% Body only: % =============================================================================
% E4 -- Workload-aware configuration selection (Tanish)
%
% Body only: % =============================================================================
% E4 -- Workload-aware configuration selection (Tanish)
%
% Body only: \input{e4_section} from the team's main file.
% Needs in the preamble: graphicx, booktabs, amsmath, cite, url.
% Figures are read from \EFourFigDir (default: figures/). Redefine it before the
% \input if the figures live elsewhere, e.g. \newcommand{\EFourFigDir}{e4/figures}.
% References are in e4_refs.bib.
% =============================================================================
\providecommand{\EFourFigDir}{figures}
\providecommand{\cfg}[1]{\texttt{#1}}

\section{Workload-Aware Configuration Selection}
\label{sec:e4}

The previous experiments establish that Torch-TensorRT latency depends on how
the graph is partitioned and on the shape of the request. This experiment asks
the question that follows from both: if the best compiler configuration moves
with the workload, can a cheap model learn where it moves, so that a deployment
picks its configuration from the batch size and sequence length instead of
committing to one setting for everything? Rice's algorithm selection
problem~\cite{rice1976algorithm} gives the framing, and cost-model-driven
autotuners such as AutoTVM~\cite{chen2018autotvm} show that the idea pays off
when the configuration space is large and the optimum is workload dependent.

Our answer for DistilBERT on an A100 is negative, and we report it as such. A
single tuned configuration is within 0.5\% of the best tested configuration on
every held-out shape. The two selectors we built reproduce that configuration
almost everywhere and, where they deviate, they mostly lose. What does matter
is tuning at all: the tuned configuration is 1.11$\times$ faster than default
Torch-TensorRT on the unmodified model and 1.6$\times$ to 2.3$\times$ faster
once part of the graph falls back to PyTorch. Getting to a result we trust took
three iterations, and two of them changed the experiment more than the
selector did. We describe them in the order they happened, because the
intermediate findings (a model with nothing to partition, a cluster that
cancels compile-heavy jobs, and a 20\% timing artifact) are useful to the rest
of the project on their own.

% -----------------------------------------------------------------------------
\subsection{Problem Formulation}
\label{sec:e4-problem}

A workload is a request shape $w=(b,s)$, with batch size $b$ and sequence
length $s$ in tokens. A configuration $c \in \mathcal{C}$ is one way of running
the model: PyTorch eager, or Torch-TensorRT with a particular set of compiler
and runtime settings. $L(c,w)$ is the warmed-up median latency of $c$ on $w$. A
selector is a function $\pi:\,w \mapsto c$ fitted on measurements from a set of
training shapes $W_\mathrm{train}$ and judged on shapes it has not seen. For a
test set $W$ we report the \emph{regret} of $\pi$ against the best tested
configuration,
\begin{equation}
  R(\pi) \;=\;
  \Bigl(\,\prod_{w\in W}
     \frac{L(\pi(w),w)}{\min_{c\in\mathcal{C}} L(c,w)}\Bigr)^{1/|W|} - 1 ,
  \label{eq:regret}
\end{equation}
that is, the geometric-mean slowdown relative to an oracle that tries every
configuration on every shape. A geometric mean keeps a 5\% loss at batch~1 and
a 5\% loss at batch~16 on equal footing, although the second is ten times
larger in milliseconds.

We compare six methods, following the evaluation plan of the proposal:
\begin{itemize}
  \item \textbf{PyTorch eager}, the uncompiled FP16 model;
  \item \textbf{Default Torch-TensorRT}, the compiler's default settings;
  \item \textbf{One fixed tuned configuration}, the configuration with the
        lowest geometric-mean slowdown over $W_\mathrm{train}$, used for every
        shape;
  \item \textbf{Lookup selector}, which returns the best configuration of the
        nearest training shape, with distance measured in
        $(\log_2 b,\log_2 s)$;
  \item \textbf{Regression selector}, which fits one latency model per
        configuration and returns the configuration with the lowest
        prediction (Section~\ref{sec:e4-selectors});
  \item \textbf{Best tested configuration}, the oracle in~\eqref{eq:regret}.
\end{itemize}
The fixed tuned configuration is the baseline that decides the experiment. A
selector is only worth deploying if it beats the best single configuration,
not merely the default one.

% -----------------------------------------------------------------------------
\subsection{Experimental Setup}
\label{sec:e4-setup}

All measurements come from one NVIDIA A100-SXM4-40GB node of the NYU HPC cloud
bursting cluster, inside a Singularity container with PyTorch~2.14.1
(CUDA~13.0)~\cite{paszke2019pytorch}, Torch-TensorRT~2.14.0~\cite{torchtrt} and
TensorRT~11.1.0~\cite{tensorrt}. The model is
\cfg{distilbert-base-uncased}~\cite{sanh2019distilbert} in FP16, wrapped so
that it takes token ids and an attention mask and returns the last hidden
state. Inputs are random token ids from a fixed seed with a full attention
mask, and every configuration receives the same tensors. Torch-TensorRT uses
its Dynamo front end, and each engine is built for one static shape, so a
(configuration, shape) pair is one compilation.

The training grid is $b\in\{1,4,8,16\}\times s\in\{128,256,512\}$, twelve
shapes. The proposal also lists 1024 tokens, but DistilBERT has 512 position
embeddings, so longer inputs cannot be run without changing the model. Eight
further shapes lie off the grid and are never used for fitting:
$(1,384)$, $(2,192)$, $(2,384)$, $(4,320)$, $(8,192)$, $(12,192)$, $(12,384)$
and $(16,448)$. None of their sequence lengths occurs in the grid, and four
of them also use a batch size (2 or 12) that does not.

Latency is measured with the project's shared benchmark routine: 20 warm-up
calls followed by 100 calls timed with CUDA events, from which we take the
median and the 95th percentile. Around it, our sweep harness records the
compile time, the time of the first call, the number of TensorRT and PyTorch
blocks in the compiled module, and the agreement of the output with eager
PyTorch (cosine similarity and largest absolute difference). The harness
appends one row per timing repeat to a CSV file and skips pairs that are
already present, so an interrupted job resumes where it stopped. That property
turned out to be necessary (Section~\ref{sec:e4-iter2}).

\begin{table}[t]
\caption{Chronology of cluster runs for this experiment (all on 12 Nov.\ 2026)}
\label{tab:e4-runs}
\centering
\footnotesize
\begin{tabular}{@{}llp{4.55cm}@{}}
\toprule
Job & Purpose & Outcome \\
\midrule
29499 & Environment & Cancelled: \cfg{--fakeroot} unpacked the 8.5\,GB image to
        network storage \\
29507 & Environment & Held by the scheduler; resubmitted \\
29509 & Environment & Packages installed (plain writable overlay) \\
29537 & Partition check & One engine, no fallback; LayerNorm exclusion had no
        effect (wrong operator name) \\
29550 & Partition check & Corrected; Table~\ref{tab:e4-partition} \\
29551 & Sweep 1 & GPU node never booted; cancelled after 12\,min \\
29586 & Sweep 1 & 83 of 340 runs, then cancelled by the cluster's
        low-utilization watchdog at 30\,min \\
29727--8 & Sweep 2 & Two more node boot failures (13\,min each) \\
29729 & Sweep 2 & 340 of 340 runs in 2\,h\,31\,min; all results below \\
\bottomrule
\end{tabular}
\end{table}

Table~\ref{tab:e4-runs} lists every cluster job behind this section. Only the
last one contributes numbers to the results; the others are the reason the
last one looks the way it does.

% -----------------------------------------------------------------------------
\subsection{Iteration 1: Is There Anything to Tune?}
\label{sec:e4-iter1}

Our plan assumed that the partitioning controls of Torch-TensorRT
(\cfg{min\_block\_size}, the choice between the fast and the global
partitioner, and operator exclusion) would produce visibly different graphs
for DistilBERT, and that the selector would arbitrate between them. Before
spending GPU hours on engines we checked that assumption with a dry-run
compile of every candidate at $(4,128)$, which reports the partition without
building anything.

The assumption was wrong. After lowering, DistilBERT consists of 127
operators (36 linear layers, 24 reshapes, 24 permutes, 13 additions, 13 layer
normalizations, six attention and six GELU operators, and a handful of
embedding and mask operations), and TensorRT supports all of them. The default
compile is therefore a single engine with no PyTorch fallback, and
\cfg{min\_block\_size} and the partitioner have nothing to act on: we verified
that \cfg{min\_block\_size}$\,=1$ and the global partitioner return the
identical one-engine graph, and removed both from the sweep.

That left two options: switch to a model with unsupported operators, which
would have cut this experiment off from the shared harness and from the other
three experiments, or keep the model and create the partitioning problem
deliberately. We chose the second and say so plainly, because it limits what
the results mean (Section~\ref{sec:e4-threats}). Torch-TensorRT lets the user
force named operators to run in PyTorch. We define three \emph{scenarios}:
\begin{itemize}
  \item \textbf{Unmodified model}: nothing excluded, one engine.
  \item \textbf{GELU in PyTorch}: the six GELU operators are excluded, which
        splits the encoder between its layers.
  \item \textbf{LayerNorm in PyTorch}: the thirteen layer normalizations are
        excluded, which cuts every layer twice.
\end{itemize}
Each scenario is evaluated separately, with its own default, its own tuned
configuration and its own selectors, as if it were a different model.

\begin{table}[t]
\caption{Partition structure of each configuration (dry-run compile)}
\label{tab:e4-partition}
\centering
\footnotesize
\setlength{\tabcolsep}{3.5pt}
\begin{tabular}{@{}lccl@{}}
\toprule
 & \multicolumn{2}{c}{Blocks} & \\
\cmidrule(lr){2-3}
Configuration & TRT & PyT. & Operators per TRT block \\
\midrule
\cfg{trt\_default}, \cfg{trt\_opt*} & 1 & 0 & 127 \\
\midrule
\cfg{gelu\_mbs5}, \cfg{gelu\_global} & 6 & 6 & 23, then $5\times19$ \\
\cfg{gelu\_mbs1} & 7 & 6 & 23, $5\times19$, 3 \\
\cfg{gelu\_mbs22} & 1 & 1 & 23 (104 ops in PyTorch) \\
\midrule
\cfg{ln\_mbs5} & 7 & 7 & 6, then $6\times14$ \\
\cfg{ln\_mbs1} & 13 & 13 & 6, then $6\times(14,\,4)$ \\
\cfg{ln\_mbs10} & 6 & 7 & $6\times14$ \\
\cfg{ln\_mbs20} & 0 & 1 & none (127 ops in PyTorch) \\
\cfg{ln\_global} & 6 & 7 & 18, then $5\times14$ \\
\bottomrule
\end{tabular}
\end{table}

Table~\ref{tab:e4-partition} shows the resulting partitions. The dry run also
pruned the configuration list a second time. With GELU excluded, the TensorRT
blocks hold 19 to 23 operators, so \cfg{min\_block\_size} values of 5, 10 and
20 all give the same graph; we kept 5 and replaced the other two by 22, the
smallest value that changes anything. Two mistakes surfaced in the same pass.
Our first LayerNorm scenario excluded \cfg{aten.layer\_norm}, but the operator
that exists after lowering is \cfg{aten.native\_layer\_norm}; the exclusion
silently matched nothing and the ``broken'' model was still one engine. And
our block counter reported zero PyTorch blocks for the global partitioner,
because that partitioner leaves fallback operators inline in the top-level
graph instead of wrapping them in submodules. Both were fixed before any
latency was measured. Without the dry run, a third of the sweep would have
been mislabeled.

\begin{table}[t]
\caption{Configurations in the sweep. Default settings are
\cfg{min\_block\_size}$\,=5$, fast partitioner, optimization level~3.}
\label{tab:e4-configs}
\centering
\footnotesize
\setlength{\tabcolsep}{3.5pt}
\begin{tabular}{@{}lll@{}}
\toprule
Scenario & Configuration & Change from scenario default \\
\midrule
all & \cfg{eager} & PyTorch, no compilation \\
\midrule
Unmodified & \cfg{trt\_default} & (default) \\
 & \cfg{trt\_opt1} & optimization level 1 \\
 & \cfg{trt\_opt5} & optimization level 5 \\
 & \cfg{trt\_cudagraphs} & CUDA Graphs at run time \\
 & \cfg{trt\_default\_rebuild} & none (second build) \\
\midrule
GELU in & \cfg{gelu\_mbs5} & (default) \\
PyTorch & \cfg{gelu\_mbs1} & \cfg{min\_block\_size} 1 \\
 & \cfg{gelu\_mbs22} & \cfg{min\_block\_size} 22 \\
 & \cfg{gelu\_global} & global partitioner \\
 & \cfg{gelu\_cudagraphs} & CUDA Graphs at run time \\
\midrule
LayerNorm & \cfg{ln\_mbs5} & (default) \\
in PyTorch & \cfg{ln\_mbs1} & \cfg{min\_block\_size} 1 \\
 & \cfg{ln\_mbs10}, \cfg{ln\_mbs20} & \cfg{min\_block\_size} 10, 20 \\
 & \cfg{ln\_global} & global partitioner \\
 & \cfg{ln\_cudagraphs} & CUDA Graphs at run time \\
\bottomrule
\end{tabular}
\end{table}

The final space has 17 configurations (Table~\ref{tab:e4-configs}). Because
partitioning offers nothing on the unmodified model, that scenario varies the
builder's optimization level instead. All three scenarios include whole-graph
CUDA Graphs~\cite{gray2019cudagraphs}, a runtime option that records the
module's stream of GPU kernel launches once and replays it as a single
submission. One more configuration, \cfg{trt\_default\_rebuild}, is the
default built a second time. It belongs to no scenario and exists to measure
how much two builds of the same configuration differ.

% -----------------------------------------------------------------------------
\subsection{Iteration 2: A First Sweep That Could Not Be Used}
\label{sec:e4-iter2}

The first sweep timed each compiled module three times with the shared
protocol. It was cancelled after 30 minutes with 83 of 340 runs done, and when
we looked at those 83 runs they were not usable either. The two problems are
unrelated.

\subsubsection{Compile-heavy jobs are cancelled}
The accounting record showed the job cancelled by the administrator account
with no message. The cause is a watchdog on the cluster that cancels a GPU job
once it is about 30 minutes old if its mean GPU utilization is 5\% or less. A
configuration sweep is exactly such a job: more than 85\% of its wall time is
TensorRT building engines, which is mostly CPU work, and timing a sub-millisecond engine
for 100 iterations keeps the GPU busy for a tenth of a second. We did not add
artificial load to stay above the threshold. The harness was already
resumable, so the second sweep was submitted as a chain of jobs, each starting
when the previous one ends and continuing from the same file.

\subsubsection{A 20\% timing artifact}
\label{sec:e4-warmup}

\begin{figure}[t]
\centering
\includegraphics[width=\columnwidth]{\EFourFigDir/fig_warmup.pdf}
\caption{Each timing repeat relative to the fastest repeat of the same
compiled module, for the five shapes both sweeps cover. (a)~With the shared
20-iteration warm-up, GPU-bound configurations start about 20\% slow and drop
to their final level one or two repeats later. (b)~With two seconds of
sustained load before timing, all repeats agree.}
\label{fig:e4-warmup}
\end{figure}

For 30 of the 79 successful runs, the three repeats of the \emph{same engine}
disagreed by more than 8\%. The disagreement has a clear shape
(Fig.~\ref{fig:e4-warmup}a): the first one or two repeats are 20\% to 25\%
slower, then the latency drops and stays down. It never goes the other way. It
affects only configurations whose latency is set by the GPU, namely the single
engine and everything run under CUDA Graphs; eager PyTorch and the
configurations with PyTorch fallback blocks, whose latency is set by the CPU
launching kernels, are stable to within 4\%.

Our reading is that the GPU needs a few hundred milliseconds of sustained
load to reach full speed after sitting idle through a compile, and that the
20-iteration warm-up is over long before that when an iteration takes
0.6\,ms. We did not confirm the mechanism (clock or power state) directly, so
we treat it as an observation about this node. The effect on a selector study
is severe either way. The differences we are trying to resolve on the
unmodified model are of the order of 10\%, and with three repeats the median
lands on the slow or the fast level essentially at random. Under these
conditions the selector would have been fitted to noise, and it would have
looked as if the best configuration changed from shape to shape. That
steady-state measurement is easy to get wrong without doing anything obviously
wrong is well documented~\cite{georges2007statistically,mytkowicz2009wrong};
this is one more instance.

% -----------------------------------------------------------------------------
\subsection{Iteration 3: The Corrected Sweep}
\label{sec:e4-iter3}

The second sweep differs from the first in three ways. Before timing, the
harness now keeps the GPU busy with the workload itself for two seconds; the
shared benchmark routine then runs unchanged, including its own 20 warm-up
calls. Each module is timed five times instead of three, and the latency we
use is the median of the five medians. And every row records the node and job
it came from, in case a chained sweep ends up spread over several machines.
The first sweep's data was discarded.

After two more failed node boots, the third job of the chain ran all 340 pairs
in one piece, in 2\,h\,31\,min on a single node. The warm-up had an unplanned
side effect: it lifted the job's GPU utilization above the watchdog threshold,
so the chain was never needed.

The data is clean. Of the 340 pairs, 320 succeeded. The 20 failures are all
\cfg{ln\_global}, which fails to build on every shape with an internal error
in Torch-TensorRT 2.14 (a missing \cfg{\_frozen\_param0} attribute on the
partitioned graph); it is excluded from selection. Across the five repeats the
median spread is 0.7\%, the 90th percentile 2.4\% and the maximum 5.8\%
(Fig.~\ref{fig:e4-warmup}b). The default configuration and its rebuild differ
by 0.3\% in the median. We therefore treat differences below about 1\% as
ties. Every compiled module agrees with eager PyTorch to a cosine similarity
of at least 0.99999, with a largest absolute difference of 0.02 in FP16.

% -----------------------------------------------------------------------------
\subsection{Selectors}
\label{sec:e4-selectors}

Both selectors are deliberately small. With twelve training shapes there is no
room for anything else.

The \emph{lookup selector} stores, for each training shape, the configuration
that was fastest there, and answers a query with the entry of the nearest
training shape in $(\log_2 b, \log_2 s)$. It has no parameters and cannot
extrapolate.

The \emph{regression selector} fits, for each configuration $c$, the latency
model
\begin{equation}
  \hat{L}_c(b,s) \;=\; \theta_{c,0} + \theta_{c,1}\,b + \theta_{c,2}\,bs
                      + \theta_{c,3}\,bs^2 ,
  \qquad \theta_c \ge 0 ,
  \label{eq:model}
\end{equation}
whose terms stand for a fixed per-call overhead, a per-sequence cost, a
per-token cost and the quadratic cost of self-attention. The coefficients are
obtained by non-negative least squares~\cite{lawson1995solving} on the
\emph{relative} error, $\sum_{w}(\hat{L}_c(w)/L(c,w)-1)^2$, so that small and
large shapes carry equal weight. Non-negativity keeps the model physically
meaningful and prevents it from predicting negative latencies outside the
grid. The selector returns $\arg\min_c \hat{L}_c(b,s)$.

We evaluate in two ways. \emph{Held-out}: fit on the twelve grid shapes, test
on the eight off-grid shapes. \emph{Leave-one-out}: for each grid shape, fit
on the other eleven and test on the one left out. To relate quality to tuning
effort, we also refit every method on subsets of $k=1,\dots,12$ grid shapes
(all subsets where there are at most 40, otherwise 40 random ones) and test on
the held-out shapes.

% -----------------------------------------------------------------------------
\subsection{Results}
\label{sec:e4-results}

\subsubsection{The latency landscape}

\begin{figure*}[t]
\centering
\includegraphics[width=\textwidth]{\EFourFigDir/fig_latency.pdf}
\caption{Median latency against batch size at 512 tokens. Eager PyTorch is
flat until the largest batch, which shows that it is limited by per-call
overhead and not by GPU work. A partitioned model inherits part of that
overhead (circles, squares); CUDA Graphs removes it (triangles). All curves meet at
batch 16, where GPU work dominates.}
\label{fig:e4-latency}
\end{figure*}

Fig.~\ref{fig:e4-latency} shows the structure that everything else follows
from. Eager PyTorch takes between 3.2 and 4.1\,ms on every shape below
$(16,448)$, almost independent of size: at these sizes an FP16 DistilBERT
forward pass on an A100 is bounded by the time it takes Python and the
framework to issue a long sequence of small kernels, not by the kernels
themselves. The single
TensorRT engine removes that overhead and is 2.68$\times$ faster than eager in
the geometric mean over all 20 shapes, from 6.6$\times$ at $(1,128)$ down to
1.07$\times$ at $(16,512)$, where the GPU is finally the bottleneck.

A partitioned model sits in between, and where it sits depends on how many
times execution returns to PyTorch. With GELU excluded (six fallback blocks)
the default configuration is 1.63$\times$ slower than the single engine; with
LayerNorm excluded (seven) it is 2.49$\times$ slower and only 7\% faster than
not compiling at all. Table~\ref{tab:e4-knobs} summarizes what each control is
worth.

\begin{table}[t]
\caption{Effect of each control on latency (geometric mean over all 20 shapes,
range over shapes in parentheses)}
\label{tab:e4-knobs}
\centering
\footnotesize
\begin{tabular}{@{}p{4.35cm}l@{}}
\toprule
Change & Latency \\
\midrule
\multicolumn{2}{@{}l}{\emph{Unmodified model, relative to default}} \\
Optimization level 5 & 1.00$\times$ (0.97--1.02) \\
Optimization level 1 & 1.06$\times$ slower (1.03--1.08) \\
CUDA Graphs & 1.11$\times$ faster (1.02--1.20) \\
\midrule
\multicolumn{2}{@{}l}{\emph{GELU in PyTorch, relative to its default}} \\
\cfg{min\_block\_size} 1 (7 engines) & 1.01$\times$ faster (0.99--1.04) \\
Global partitioner & 1.02$\times$ slower (0.98--1.05) \\
\cfg{min\_block\_size} 22 (1 engine) & 1.74$\times$ slower (1.07--2.25) \\
CUDA Graphs & 1.63$\times$ faster (1.06--3.27) \\
\midrule
\multicolumn{2}{@{}l}{\emph{LayerNorm in PyTorch, relative to its default}} \\
\cfg{min\_block\_size} 1 (13 engines) & 1.05$\times$ faster (0.99--1.11) \\
\cfg{min\_block\_size} 10 (6 engines) & 1.03$\times$ slower (1.00--1.05) \\
\cfg{min\_block\_size} 20 (no engine) & 1.28$\times$ slower (1.04--1.47) \\
CUDA Graphs & 2.26$\times$ faster (1.08--5.02) \\
\bottomrule
\end{tabular}
\end{table}

Two things stand out. First, the partitioning controls barely matter as long
as the engines survive: \cfg{min\_block\_size} and the partitioner move
latency by 5\% or less, and their only large effect is the destructive one,
when a high threshold discards engines and pushes the model back towards
eager. Second, one runtime option dwarfs all of them. CUDA Graphs makes the
GELU-partitioned model 1.63$\times$ faster and the LayerNorm-partitioned one
2.26$\times$ faster, to within 11\% and 23\% of the single engine under the
same option. This is consistent with the picture above: what a fallback block
costs is launch overhead, and replaying a recorded launch sequence removes it
no matter how many blocks there are.

\begin{figure*}[t]
\centering
\includegraphics[width=\textwidth]{\EFourFigDir/fig_gain.pdf}
\caption{Speedup of the best tested configuration over the scenario's default
on the training grid. Tuning is worth most at small shapes and when the graph
is partitioned; it fades as the GPU becomes the bottleneck.}
\label{fig:e4-gain}
\end{figure*}

Fig.~\ref{fig:e4-gain} shows the same thing from the point of view of tuning.
The \emph{size} of the gain is strongly workload dependent, from 5.0$\times$
at $(1,128)$ with LayerNorm excluded down to a few percent at $(16,512)$. The
\emph{identity} of the best configuration is not, and only the latter is
something a selector can exploit.

\subsubsection{Selection on held-out shapes}

\begin{table}[t]
\caption{Held-out shapes: regret $R$ from~\eqref{eq:regret} and speedup over
default Torch-TensorRT. GELU and LayerNorm denote the scenarios in which that
operator runs in PyTorch.}
\label{tab:e4-heldout}
\centering
\footnotesize
\setlength{\tabcolsep}{2.6pt}
\begin{tabular}{@{}lrrrrrr@{}}
\toprule
 & \multicolumn{2}{c}{Unmodified} & \multicolumn{2}{c}{GELU}
 & \multicolumn{2}{c}{LayerNorm} \\
\cmidrule(lr){2-3}\cmidrule(lr){4-5}\cmidrule(l){6-7}
Method & $R$\,(\%) & speedup & $R$\,(\%) & speedup & $R$\,(\%) & speedup \\
\midrule
PyTorch eager        & 209.8 & 0.36$\times$ & 180.9 & 0.57$\times$ & 153.8 & 0.92$\times$ \\
Default TRT          &  11.0 & 1.00$\times$ &  61.3 & 1.00$\times$ & 132.7 & 1.00$\times$ \\
Fixed tuned          &   0.0 & 1.11$\times$ &   0.0 & 1.61$\times$ &   0.5 & 2.32$\times$ \\
Lookup sel.          &   0.0 & 1.11$\times$ &   3.2 & 1.56$\times$ &   1.4 & 2.29$\times$ \\
Regression sel.      &   1.5 & 1.09$\times$ &   0.5 & 1.60$\times$ &   0.0 & 2.33$\times$ \\
Best tested          &   0.0 & 1.11$\times$ &   0.0 & 1.61$\times$ &   0.0 & 2.33$\times$ \\
\bottomrule
\end{tabular}
\end{table}

\begin{figure}[t]
\centering
\includegraphics[width=\columnwidth]{\EFourFigDir/fig_regret.pdf}
\caption{Regret on the eight held-out shapes for the three methods that use
tuning data. The dashed line is the timing noise (0.7\%). No selector beats
the fixed tuned configuration by more than the noise.}
\label{fig:e4-regret}
\end{figure}

Table~\ref{tab:e4-heldout} and Fig.~\ref{fig:e4-regret} give the main result.
In all three scenarios the fixed tuned configuration is the CUDA Graphs
variant, and it is the best tested configuration on 23 of the 24
scenario--shape combinations. Its regret is 0.0\%, 0.0\% and 0.5\%. The
selectors range from 0.0\% to 3.2\%; relative to the fixed tuned configuration
they achieve between 0.97$\times$ and 1.01$\times$. The single case in which a
selector is ahead (regression, LayerNorm scenario, by 0.5\%) is below the
noise floor. Leave-one-out evaluation on the grid agrees: the fixed tuned
configuration has a regret of 0.0\%, 0.1\% and 0.9\%, the lookup selector
makes identical choices, and the regression selector scores 0.8\%, 1.0\% and
0.0\%.

Against default Torch-TensorRT, on the other hand, the gain is real and
sizeable: 1.11$\times$ on the unmodified model, 1.61$\times$ with GELU
excluded and 2.32$\times$ with LayerNorm excluded.

\subsubsection{Why there is nothing to select}

\begin{figure}[t]
\centering
\includegraphics[width=\columnwidth]{\EFourFigDir/fig_choice.pdf}
\caption{Fastest configuration at each of the 20 shapes in the two
partitioned scenarios. The only shape dependence is a switch to eager PyTorch
at the largest shapes. Circles mark held-out shapes where the lookup selector
chose differently from the best configuration, with the resulting loss.}
\label{fig:e4-choice}
\end{figure}

On the unmodified model, \cfg{trt\_cudagraphs} is the fastest configuration on
all 20 shapes. In the partitioned scenarios the CUDA Graphs variant is fastest
everywhere except at the top-right corner of the shape space
(Fig.~\ref{fig:e4-choice}): at $(16,512)$ in both scenarios and at $(16,448)$
with LayerNorm excluded, eager PyTorch is faster by 1\% to 11\%, because at
that size the per-call overhead no longer matters and the compiled module
pays for crossing between engines and PyTorch.

That one crossover is the entire workload dependence in our data, and the
selectors handle it poorly. Across the three scenarios the two selectors chose
eager for a held-out shape seven times, and it was the right call twice. The
lookup selector assigns $(12,384)$ to its nearest grid neighbor $(16,512)$ and
therefore to eager, at a cost of 23\% on that shape with GELU excluded and
12\% with LayerNorm excluded. The regression selector does better on the
partitioned scenarios but sends $(16,448)$ to eager on the unmodified model,
where the engine is 13\% faster. The underlying problem is that the boundary
is defined by a single training point at the edge of the grid, and that the
latency model is least accurate for exactly the configuration on the other
side of it: its mean absolute error on held-out shapes is 3\% to 4\% for the
GPU-bound configurations but 13\% for eager, whose flat-then-rising curve the
four terms of~\eqref{eq:model} do not fit well.

\subsubsection{Tuning cost}

\begin{figure*}[t]
\centering
\includegraphics[width=\textwidth]{\EFourFigDir/fig_budget.pdf}
\caption{Held-out regret when only $k$ of the twelve grid shapes are measured
(mean over subsets). Two shapes are enough to identify the fixed tuned
configuration. With one shape the outcome depends on which: if it is
$(16,512)$, every method selects eager PyTorch.}
\label{fig:e4-budget}
\end{figure*}

\begin{table}[t]
\caption{Tuning cost and break-even against default Torch-TensorRT, per
scenario}
\label{tab:e4-cost}
\centering
\footnotesize
\setlength{\tabcolsep}{4pt}
\begin{tabular}{@{}lrrr@{}}
\toprule
 & Unmodified & GELU & LayerNorm \\
\midrule
Saving per request (ms)        & 0.12 & 0.70 & 1.76 \\
\midrule
Full grid, 12 shapes (GPU-min) & 16.6 & 30.1 & 33.3 \\
\quad break-even (requests)    & $8.3\times10^6$ & $2.6\times10^6$ & $1.1\times10^6$ \\
\midrule
Two shapes (GPU-min)           & 2.7 & 5.1 & 5.5 \\
\quad break-even (requests)    & $1.4\times10^6$ & $4.3\times10^5$ & $1.9\times10^5$ \\
\bottomrule
\end{tabular}
\end{table}

The cost of a measurement is the engine build, the first call, the warm-up
and the timing. Measuring every configuration of a scenario on all twelve
grid shapes costs 17 to 33 GPU-minutes (Table~\ref{tab:e4-cost}), of which
engine builds are the bulk: a single-engine build takes 8 to 14\,s at the
default optimization level and 32\,s at level~5, a partitioned one about
30\,s, and the 13-engine \cfg{ln\_mbs1} about 57\,s.
Fig.~\ref{fig:e4-budget} shows that most of this effort is unnecessary. Two
measured shapes identify the fixed tuned configuration in every scenario, and
the ten further shapes buy nothing for it; the selectors need more data and
still end up level or behind. Dividing the cost by the mean saving per request
on the held-out shapes gives the number of requests after which tuning has
paid for itself: about 1.4 million on the unmodified model, where the saving
is 0.12\,ms, and 0.2 to 0.4 million in the partitioned scenarios.

% -----------------------------------------------------------------------------
\subsection{Discussion}
\label{sec:e4-discussion}

The proposal's hypothesis for this experiment was that a lookup table or a
regression model would recover most of the gap between a fixed configuration
and the best one. For this model the hypothesis fails for an uninteresting
reason: the gap is not there. We think the reason generalizes somewhat beyond
our setup, and it is worth stating.

A selector needs configurations whose latency curves \emph{cross}. In our
space, the configurations differ in how much per-call overhead they carry, and
overhead is a constant offset: a configuration with less of it is faster at
every shape, by an amount that shrinks as the shape grows but does not change
sign. CUDA Graphs removes the most overhead and therefore wins everywhere it
is applicable. The only crossing we observe involves eager PyTorch at the
largest shapes, and it is shallow (at most 11\%) and located at the edge of
the measured region. Per-workload selection becomes interesting when
configurations trade a fixed cost against a marginal one, for instance
precision modes with different kernel speeds, engines built with dynamic
versus static shape profiles, or models where a fallback block does real
computation instead of just interrupting the engine. None of these is in our
space.

For the project as a whole, three findings from this experiment matter more
than the selector:
\begin{enumerate}
  \item DistilBERT converts to a single TensorRT engine. Any partitioning
        effect studied on it has to be induced.
  \item On this model the cost of a partition boundary is launch overhead,
        and CUDA Graphs recovers most of it. Whether a graph is ``well
        partitioned'' matters far less with that option enabled.
  \item The shared benchmark's 20-iteration warm-up is too short for
        sub-millisecond engines on this hardware and overstates their latency
        by about 20\% in some repeats. Experiments that compare a single
        engine against PyTorch with it will understate the speedup.
\end{enumerate}

% -----------------------------------------------------------------------------
\subsection{Threats to Validity}
\label{sec:e4-threats}

\emph{The partitioned scenarios are artificial.} GELU and LayerNorm convert
without difficulty, and nobody would exclude them in practice. The scenarios
reproduce the structure of a fallback (execution leaves the engine and comes
back) but not necessarily its cost profile, since a genuinely unsupported
operator may be expensive in itself.

\emph{One model, one GPU, one precision.} Everything here is DistilBERT in
FP16 on one A100 node, with sequences of at most 512 tokens. All 340 pairs
were measured in a single job on a single machine, which removes
machine-to-machine variation from the comparison but also means we have not
observed it.

\emph{Static shapes.} Every engine is built for exactly one shape. A
deployment that serves many shapes from one engine with a dynamic profile
faces a different trade-off, which this experiment does not measure.

\emph{Drift in the CPU-bound baselines.} The held-out shapes were measured
in the last hour of the job, and the CPU-bound configurations look slower
there than on the grid. Relative to the grid-fitted model~\eqref{eq:model},
eager PyTorch is 15\% slower on the held-out shapes on average, the
LayerNorm-partitioned configurations 8\% to 11\%, and the GPU-bound
configurations 1\%. Some of this is model error, but neighboring shapes point
the same way (eager takes 3.25\,ms at $(8,256)$ on the grid and 4.06\,ms at
$(8,192)$ held out), so host-side performance probably degraded during the
job. Two consequences follow. The held-out speedups over eager and over the
partitioned defaults are somewhat inflated; the leave-one-out figures on the
grid (2.22$\times$ over the default for LayerNorm, against 2.32$\times$ held
out) are the more conservative ones. And the position of the eager crossover
on the held-out shapes is uncertain: if eager were in fact 15\% faster at
$(16,448)$, it would be the best choice there in all three scenarios, the
fixed tuned configuration would lose 2\% to 20\% on that one shape, and a
selector choosing eager there would be right. This would move the
held-out regret of the fixed tuned configuration to at most 2.3\%. It would not
change the picture at the other seven shapes or on the grid.

\emph{Build times share a cache.} TensorRT reuses layer timings within a
process, so the first build of the job took 40\,s and later builds of the same
graph 6 to 13\,s. The tuning costs in Table~\ref{tab:e4-cost} are therefore
optimistic for a cold start.

\emph{Memory.} The recorded peak memory counts PyTorch's allocator only and
excludes memory held by TensorRT engines, so we do not report a memory
comparison.

% -----------------------------------------------------------------------------
\subsection{Summary}
\label{sec:e4-summary}

We built a resumable sweep harness, two workload-aware selectors and an
evaluation on held-out shapes, and measured 17 configurations on 20 shapes.
For DistilBERT on an A100, tuning the configuration once is worth
1.11$\times$ to 2.32$\times$ over default Torch-TensorRT, two measured shapes
are enough to do it, and selecting the configuration per workload adds nothing
beyond that. The result rests on a measurement protocol that had to be
corrected first; with the original warm-up the same experiment would have
reported a workload-dependent optimum that does not exist.
 from the team's main file.
% Needs in the preamble: graphicx, booktabs, amsmath, cite, url.
% Figures are read from \EFourFigDir (default: figures/). Redefine it before the
% \input if the figures live elsewhere, e.g. \newcommand{\EFourFigDir}{e4/figures}.
% References are in e4_refs.bib.
% =============================================================================
\providecommand{\EFourFigDir}{figures}
\providecommand{\cfg}[1]{\texttt{#1}}

\section{Workload-Aware Configuration Selection}
\label{sec:e4}

The previous experiments establish that Torch-TensorRT latency depends on how
the graph is partitioned and on the shape of the request. This experiment asks
the question that follows from both: if the best compiler configuration moves
with the workload, can a cheap model learn where it moves, so that a deployment
picks its configuration from the batch size and sequence length instead of
committing to one setting for everything? Rice's algorithm selection
problem~\cite{rice1976algorithm} gives the framing, and cost-model-driven
autotuners such as AutoTVM~\cite{chen2018autotvm} show that the idea pays off
when the configuration space is large and the optimum is workload dependent.

Our answer for DistilBERT on an A100 is negative, and we report it as such. A
single tuned configuration is within 0.5\% of the best tested configuration on
every held-out shape. The two selectors we built reproduce that configuration
almost everywhere and, where they deviate, they mostly lose. What does matter
is tuning at all: the tuned configuration is 1.11$\times$ faster than default
Torch-TensorRT on the unmodified model and 1.6$\times$ to 2.3$\times$ faster
once part of the graph falls back to PyTorch. Getting to a result we trust took
three iterations, and two of them changed the experiment more than the
selector did. We describe them in the order they happened, because the
intermediate findings (a model with nothing to partition, a cluster that
cancels compile-heavy jobs, and a 20\% timing artifact) are useful to the rest
of the project on their own.

% -----------------------------------------------------------------------------
\subsection{Problem Formulation}
\label{sec:e4-problem}

A workload is a request shape $w=(b,s)$, with batch size $b$ and sequence
length $s$ in tokens. A configuration $c \in \mathcal{C}$ is one way of running
the model: PyTorch eager, or Torch-TensorRT with a particular set of compiler
and runtime settings. $L(c,w)$ is the warmed-up median latency of $c$ on $w$. A
selector is a function $\pi:\,w \mapsto c$ fitted on measurements from a set of
training shapes $W_\mathrm{train}$ and judged on shapes it has not seen. For a
test set $W$ we report the \emph{regret} of $\pi$ against the best tested
configuration,
\begin{equation}
  R(\pi) \;=\;
  \Bigl(\,\prod_{w\in W}
     \frac{L(\pi(w),w)}{\min_{c\in\mathcal{C}} L(c,w)}\Bigr)^{1/|W|} - 1 ,
  \label{eq:regret}
\end{equation}
that is, the geometric-mean slowdown relative to an oracle that tries every
configuration on every shape. A geometric mean keeps a 5\% loss at batch~1 and
a 5\% loss at batch~16 on equal footing, although the second is ten times
larger in milliseconds.

We compare six methods, following the evaluation plan of the proposal:
\begin{itemize}
  \item \textbf{PyTorch eager}, the uncompiled FP16 model;
  \item \textbf{Default Torch-TensorRT}, the compiler's default settings;
  \item \textbf{One fixed tuned configuration}, the configuration with the
        lowest geometric-mean slowdown over $W_\mathrm{train}$, used for every
        shape;
  \item \textbf{Lookup selector}, which returns the best configuration of the
        nearest training shape, with distance measured in
        $(\log_2 b,\log_2 s)$;
  \item \textbf{Regression selector}, which fits one latency model per
        configuration and returns the configuration with the lowest
        prediction (Section~\ref{sec:e4-selectors});
  \item \textbf{Best tested configuration}, the oracle in~\eqref{eq:regret}.
\end{itemize}
The fixed tuned configuration is the baseline that decides the experiment. A
selector is only worth deploying if it beats the best single configuration,
not merely the default one.

% -----------------------------------------------------------------------------
\subsection{Experimental Setup}
\label{sec:e4-setup}

All measurements come from one NVIDIA A100-SXM4-40GB node of the NYU HPC cloud
bursting cluster, inside a Singularity container with PyTorch~2.14.1
(CUDA~13.0)~\cite{paszke2019pytorch}, Torch-TensorRT~2.14.0~\cite{torchtrt} and
TensorRT~11.1.0~\cite{tensorrt}. The model is
\cfg{distilbert-base-uncased}~\cite{sanh2019distilbert} in FP16, wrapped so
that it takes token ids and an attention mask and returns the last hidden
state. Inputs are random token ids from a fixed seed with a full attention
mask, and every configuration receives the same tensors. Torch-TensorRT uses
its Dynamo front end, and each engine is built for one static shape, so a
(configuration, shape) pair is one compilation.

The training grid is $b\in\{1,4,8,16\}\times s\in\{128,256,512\}$, twelve
shapes. The proposal also lists 1024 tokens, but DistilBERT has 512 position
embeddings, so longer inputs cannot be run without changing the model. Eight
further shapes lie off the grid and are never used for fitting:
$(1,384)$, $(2,192)$, $(2,384)$, $(4,320)$, $(8,192)$, $(12,192)$, $(12,384)$
and $(16,448)$. None of their sequence lengths occurs in the grid, and four
of them also use a batch size (2 or 12) that does not.

Latency is measured with the project's shared benchmark routine: 20 warm-up
calls followed by 100 calls timed with CUDA events, from which we take the
median and the 95th percentile. Around it, our sweep harness records the
compile time, the time of the first call, the number of TensorRT and PyTorch
blocks in the compiled module, and the agreement of the output with eager
PyTorch (cosine similarity and largest absolute difference). The harness
appends one row per timing repeat to a CSV file and skips pairs that are
already present, so an interrupted job resumes where it stopped. That property
turned out to be necessary (Section~\ref{sec:e4-iter2}).

\begin{table}[t]
\caption{Chronology of cluster runs for this experiment (all on 12 Nov.\ 2026)}
\label{tab:e4-runs}
\centering
\footnotesize
\begin{tabular}{@{}llp{4.55cm}@{}}
\toprule
Job & Purpose & Outcome \\
\midrule
29499 & Environment & Cancelled: \cfg{--fakeroot} unpacked the 8.5\,GB image to
        network storage \\
29507 & Environment & Held by the scheduler; resubmitted \\
29509 & Environment & Packages installed (plain writable overlay) \\
29537 & Partition check & One engine, no fallback; LayerNorm exclusion had no
        effect (wrong operator name) \\
29550 & Partition check & Corrected; Table~\ref{tab:e4-partition} \\
29551 & Sweep 1 & GPU node never booted; cancelled after 12\,min \\
29586 & Sweep 1 & 83 of 340 runs, then cancelled by the cluster's
        low-utilization watchdog at 30\,min \\
29727--8 & Sweep 2 & Two more node boot failures (13\,min each) \\
29729 & Sweep 2 & 340 of 340 runs in 2\,h\,31\,min; all results below \\
\bottomrule
\end{tabular}
\end{table}

Table~\ref{tab:e4-runs} lists every cluster job behind this section. Only the
last one contributes numbers to the results; the others are the reason the
last one looks the way it does.

% -----------------------------------------------------------------------------
\subsection{Iteration 1: Is There Anything to Tune?}
\label{sec:e4-iter1}

Our plan assumed that the partitioning controls of Torch-TensorRT
(\cfg{min\_block\_size}, the choice between the fast and the global
partitioner, and operator exclusion) would produce visibly different graphs
for DistilBERT, and that the selector would arbitrate between them. Before
spending GPU hours on engines we checked that assumption with a dry-run
compile of every candidate at $(4,128)$, which reports the partition without
building anything.

The assumption was wrong. After lowering, DistilBERT consists of 127
operators (36 linear layers, 24 reshapes, 24 permutes, 13 additions, 13 layer
normalizations, six attention and six GELU operators, and a handful of
embedding and mask operations), and TensorRT supports all of them. The default
compile is therefore a single engine with no PyTorch fallback, and
\cfg{min\_block\_size} and the partitioner have nothing to act on: we verified
that \cfg{min\_block\_size}$\,=1$ and the global partitioner return the
identical one-engine graph, and removed both from the sweep.

That left two options: switch to a model with unsupported operators, which
would have cut this experiment off from the shared harness and from the other
three experiments, or keep the model and create the partitioning problem
deliberately. We chose the second and say so plainly, because it limits what
the results mean (Section~\ref{sec:e4-threats}). Torch-TensorRT lets the user
force named operators to run in PyTorch. We define three \emph{scenarios}:
\begin{itemize}
  \item \textbf{Unmodified model}: nothing excluded, one engine.
  \item \textbf{GELU in PyTorch}: the six GELU operators are excluded, which
        splits the encoder between its layers.
  \item \textbf{LayerNorm in PyTorch}: the thirteen layer normalizations are
        excluded, which cuts every layer twice.
\end{itemize}
Each scenario is evaluated separately, with its own default, its own tuned
configuration and its own selectors, as if it were a different model.

\begin{table}[t]
\caption{Partition structure of each configuration (dry-run compile)}
\label{tab:e4-partition}
\centering
\footnotesize
\setlength{\tabcolsep}{3.5pt}
\begin{tabular}{@{}lccl@{}}
\toprule
 & \multicolumn{2}{c}{Blocks} & \\
\cmidrule(lr){2-3}
Configuration & TRT & PyT. & Operators per TRT block \\
\midrule
\cfg{trt\_default}, \cfg{trt\_opt*} & 1 & 0 & 127 \\
\midrule
\cfg{gelu\_mbs5}, \cfg{gelu\_global} & 6 & 6 & 23, then $5\times19$ \\
\cfg{gelu\_mbs1} & 7 & 6 & 23, $5\times19$, 3 \\
\cfg{gelu\_mbs22} & 1 & 1 & 23 (104 ops in PyTorch) \\
\midrule
\cfg{ln\_mbs5} & 7 & 7 & 6, then $6\times14$ \\
\cfg{ln\_mbs1} & 13 & 13 & 6, then $6\times(14,\,4)$ \\
\cfg{ln\_mbs10} & 6 & 7 & $6\times14$ \\
\cfg{ln\_mbs20} & 0 & 1 & none (127 ops in PyTorch) \\
\cfg{ln\_global} & 6 & 7 & 18, then $5\times14$ \\
\bottomrule
\end{tabular}
\end{table}

Table~\ref{tab:e4-partition} shows the resulting partitions. The dry run also
pruned the configuration list a second time. With GELU excluded, the TensorRT
blocks hold 19 to 23 operators, so \cfg{min\_block\_size} values of 5, 10 and
20 all give the same graph; we kept 5 and replaced the other two by 22, the
smallest value that changes anything. Two mistakes surfaced in the same pass.
Our first LayerNorm scenario excluded \cfg{aten.layer\_norm}, but the operator
that exists after lowering is \cfg{aten.native\_layer\_norm}; the exclusion
silently matched nothing and the ``broken'' model was still one engine. And
our block counter reported zero PyTorch blocks for the global partitioner,
because that partitioner leaves fallback operators inline in the top-level
graph instead of wrapping them in submodules. Both were fixed before any
latency was measured. Without the dry run, a third of the sweep would have
been mislabeled.

\begin{table}[t]
\caption{Configurations in the sweep. Default settings are
\cfg{min\_block\_size}$\,=5$, fast partitioner, optimization level~3.}
\label{tab:e4-configs}
\centering
\footnotesize
\setlength{\tabcolsep}{3.5pt}
\begin{tabular}{@{}lll@{}}
\toprule
Scenario & Configuration & Change from scenario default \\
\midrule
all & \cfg{eager} & PyTorch, no compilation \\
\midrule
Unmodified & \cfg{trt\_default} & (default) \\
 & \cfg{trt\_opt1} & optimization level 1 \\
 & \cfg{trt\_opt5} & optimization level 5 \\
 & \cfg{trt\_cudagraphs} & CUDA Graphs at run time \\
 & \cfg{trt\_default\_rebuild} & none (second build) \\
\midrule
GELU in & \cfg{gelu\_mbs5} & (default) \\
PyTorch & \cfg{gelu\_mbs1} & \cfg{min\_block\_size} 1 \\
 & \cfg{gelu\_mbs22} & \cfg{min\_block\_size} 22 \\
 & \cfg{gelu\_global} & global partitioner \\
 & \cfg{gelu\_cudagraphs} & CUDA Graphs at run time \\
\midrule
LayerNorm & \cfg{ln\_mbs5} & (default) \\
in PyTorch & \cfg{ln\_mbs1} & \cfg{min\_block\_size} 1 \\
 & \cfg{ln\_mbs10}, \cfg{ln\_mbs20} & \cfg{min\_block\_size} 10, 20 \\
 & \cfg{ln\_global} & global partitioner \\
 & \cfg{ln\_cudagraphs} & CUDA Graphs at run time \\
\bottomrule
\end{tabular}
\end{table}

The final space has 17 configurations (Table~\ref{tab:e4-configs}). Because
partitioning offers nothing on the unmodified model, that scenario varies the
builder's optimization level instead. All three scenarios include whole-graph
CUDA Graphs~\cite{gray2019cudagraphs}, a runtime option that records the
module's stream of GPU kernel launches once and replays it as a single
submission. One more configuration, \cfg{trt\_default\_rebuild}, is the
default built a second time. It belongs to no scenario and exists to measure
how much two builds of the same configuration differ.

% -----------------------------------------------------------------------------
\subsection{Iteration 2: A First Sweep That Could Not Be Used}
\label{sec:e4-iter2}

The first sweep timed each compiled module three times with the shared
protocol. It was cancelled after 30 minutes with 83 of 340 runs done, and when
we looked at those 83 runs they were not usable either. The two problems are
unrelated.

\subsubsection{Compile-heavy jobs are cancelled}
The accounting record showed the job cancelled by the administrator account
with no message. The cause is a watchdog on the cluster that cancels a GPU job
once it is about 30 minutes old if its mean GPU utilization is 5\% or less. A
configuration sweep is exactly such a job: more than 85\% of its wall time is
TensorRT building engines, which is mostly CPU work, and timing a sub-millisecond engine
for 100 iterations keeps the GPU busy for a tenth of a second. We did not add
artificial load to stay above the threshold. The harness was already
resumable, so the second sweep was submitted as a chain of jobs, each starting
when the previous one ends and continuing from the same file.

\subsubsection{A 20\% timing artifact}
\label{sec:e4-warmup}

\begin{figure}[t]
\centering
\includegraphics[width=\columnwidth]{\EFourFigDir/fig_warmup.pdf}
\caption{Each timing repeat relative to the fastest repeat of the same
compiled module, for the five shapes both sweeps cover. (a)~With the shared
20-iteration warm-up, GPU-bound configurations start about 20\% slow and drop
to their final level one or two repeats later. (b)~With two seconds of
sustained load before timing, all repeats agree.}
\label{fig:e4-warmup}
\end{figure}

For 30 of the 79 successful runs, the three repeats of the \emph{same engine}
disagreed by more than 8\%. The disagreement has a clear shape
(Fig.~\ref{fig:e4-warmup}a): the first one or two repeats are 20\% to 25\%
slower, then the latency drops and stays down. It never goes the other way. It
affects only configurations whose latency is set by the GPU, namely the single
engine and everything run under CUDA Graphs; eager PyTorch and the
configurations with PyTorch fallback blocks, whose latency is set by the CPU
launching kernels, are stable to within 4\%.

Our reading is that the GPU needs a few hundred milliseconds of sustained
load to reach full speed after sitting idle through a compile, and that the
20-iteration warm-up is over long before that when an iteration takes
0.6\,ms. We did not confirm the mechanism (clock or power state) directly, so
we treat it as an observation about this node. The effect on a selector study
is severe either way. The differences we are trying to resolve on the
unmodified model are of the order of 10\%, and with three repeats the median
lands on the slow or the fast level essentially at random. Under these
conditions the selector would have been fitted to noise, and it would have
looked as if the best configuration changed from shape to shape. That
steady-state measurement is easy to get wrong without doing anything obviously
wrong is well documented~\cite{georges2007statistically,mytkowicz2009wrong};
this is one more instance.

% -----------------------------------------------------------------------------
\subsection{Iteration 3: The Corrected Sweep}
\label{sec:e4-iter3}

The second sweep differs from the first in three ways. Before timing, the
harness now keeps the GPU busy with the workload itself for two seconds; the
shared benchmark routine then runs unchanged, including its own 20 warm-up
calls. Each module is timed five times instead of three, and the latency we
use is the median of the five medians. And every row records the node and job
it came from, in case a chained sweep ends up spread over several machines.
The first sweep's data was discarded.

After two more failed node boots, the third job of the chain ran all 340 pairs
in one piece, in 2\,h\,31\,min on a single node. The warm-up had an unplanned
side effect: it lifted the job's GPU utilization above the watchdog threshold,
so the chain was never needed.

The data is clean. Of the 340 pairs, 320 succeeded. The 20 failures are all
\cfg{ln\_global}, which fails to build on every shape with an internal error
in Torch-TensorRT 2.14 (a missing \cfg{\_frozen\_param0} attribute on the
partitioned graph); it is excluded from selection. Across the five repeats the
median spread is 0.7\%, the 90th percentile 2.4\% and the maximum 5.8\%
(Fig.~\ref{fig:e4-warmup}b). The default configuration and its rebuild differ
by 0.3\% in the median. We therefore treat differences below about 1\% as
ties. Every compiled module agrees with eager PyTorch to a cosine similarity
of at least 0.99999, with a largest absolute difference of 0.02 in FP16.

% -----------------------------------------------------------------------------
\subsection{Selectors}
\label{sec:e4-selectors}

Both selectors are deliberately small. With twelve training shapes there is no
room for anything else.

The \emph{lookup selector} stores, for each training shape, the configuration
that was fastest there, and answers a query with the entry of the nearest
training shape in $(\log_2 b, \log_2 s)$. It has no parameters and cannot
extrapolate.

The \emph{regression selector} fits, for each configuration $c$, the latency
model
\begin{equation}
  \hat{L}_c(b,s) \;=\; \theta_{c,0} + \theta_{c,1}\,b + \theta_{c,2}\,bs
                      + \theta_{c,3}\,bs^2 ,
  \qquad \theta_c \ge 0 ,
  \label{eq:model}
\end{equation}
whose terms stand for a fixed per-call overhead, a per-sequence cost, a
per-token cost and the quadratic cost of self-attention. The coefficients are
obtained by non-negative least squares~\cite{lawson1995solving} on the
\emph{relative} error, $\sum_{w}(\hat{L}_c(w)/L(c,w)-1)^2$, so that small and
large shapes carry equal weight. Non-negativity keeps the model physically
meaningful and prevents it from predicting negative latencies outside the
grid. The selector returns $\arg\min_c \hat{L}_c(b,s)$.

We evaluate in two ways. \emph{Held-out}: fit on the twelve grid shapes, test
on the eight off-grid shapes. \emph{Leave-one-out}: for each grid shape, fit
on the other eleven and test on the one left out. To relate quality to tuning
effort, we also refit every method on subsets of $k=1,\dots,12$ grid shapes
(all subsets where there are at most 40, otherwise 40 random ones) and test on
the held-out shapes.

% -----------------------------------------------------------------------------
\subsection{Results}
\label{sec:e4-results}

\subsubsection{The latency landscape}

\begin{figure*}[t]
\centering
\includegraphics[width=\textwidth]{\EFourFigDir/fig_latency.pdf}
\caption{Median latency against batch size at 512 tokens. Eager PyTorch is
flat until the largest batch, which shows that it is limited by per-call
overhead and not by GPU work. A partitioned model inherits part of that
overhead (circles, squares); CUDA Graphs removes it (triangles). All curves meet at
batch 16, where GPU work dominates.}
\label{fig:e4-latency}
\end{figure*}

Fig.~\ref{fig:e4-latency} shows the structure that everything else follows
from. Eager PyTorch takes between 3.2 and 4.1\,ms on every shape below
$(16,448)$, almost independent of size: at these sizes an FP16 DistilBERT
forward pass on an A100 is bounded by the time it takes Python and the
framework to issue a long sequence of small kernels, not by the kernels
themselves. The single
TensorRT engine removes that overhead and is 2.68$\times$ faster than eager in
the geometric mean over all 20 shapes, from 6.6$\times$ at $(1,128)$ down to
1.07$\times$ at $(16,512)$, where the GPU is finally the bottleneck.

A partitioned model sits in between, and where it sits depends on how many
times execution returns to PyTorch. With GELU excluded (six fallback blocks)
the default configuration is 1.63$\times$ slower than the single engine; with
LayerNorm excluded (seven) it is 2.49$\times$ slower and only 7\% faster than
not compiling at all. Table~\ref{tab:e4-knobs} summarizes what each control is
worth.

\begin{table}[t]
\caption{Effect of each control on latency (geometric mean over all 20 shapes,
range over shapes in parentheses)}
\label{tab:e4-knobs}
\centering
\footnotesize
\begin{tabular}{@{}p{4.35cm}l@{}}
\toprule
Change & Latency \\
\midrule
\multicolumn{2}{@{}l}{\emph{Unmodified model, relative to default}} \\
Optimization level 5 & 1.00$\times$ (0.97--1.02) \\
Optimization level 1 & 1.06$\times$ slower (1.03--1.08) \\
CUDA Graphs & 1.11$\times$ faster (1.02--1.20) \\
\midrule
\multicolumn{2}{@{}l}{\emph{GELU in PyTorch, relative to its default}} \\
\cfg{min\_block\_size} 1 (7 engines) & 1.01$\times$ faster (0.99--1.04) \\
Global partitioner & 1.02$\times$ slower (0.98--1.05) \\
\cfg{min\_block\_size} 22 (1 engine) & 1.74$\times$ slower (1.07--2.25) \\
CUDA Graphs & 1.63$\times$ faster (1.06--3.27) \\
\midrule
\multicolumn{2}{@{}l}{\emph{LayerNorm in PyTorch, relative to its default}} \\
\cfg{min\_block\_size} 1 (13 engines) & 1.05$\times$ faster (0.99--1.11) \\
\cfg{min\_block\_size} 10 (6 engines) & 1.03$\times$ slower (1.00--1.05) \\
\cfg{min\_block\_size} 20 (no engine) & 1.28$\times$ slower (1.04--1.47) \\
CUDA Graphs & 2.26$\times$ faster (1.08--5.02) \\
\bottomrule
\end{tabular}
\end{table}

Two things stand out. First, the partitioning controls barely matter as long
as the engines survive: \cfg{min\_block\_size} and the partitioner move
latency by 5\% or less, and their only large effect is the destructive one,
when a high threshold discards engines and pushes the model back towards
eager. Second, one runtime option dwarfs all of them. CUDA Graphs makes the
GELU-partitioned model 1.63$\times$ faster and the LayerNorm-partitioned one
2.26$\times$ faster, to within 11\% and 23\% of the single engine under the
same option. This is consistent with the picture above: what a fallback block
costs is launch overhead, and replaying a recorded launch sequence removes it
no matter how many blocks there are.

\begin{figure*}[t]
\centering
\includegraphics[width=\textwidth]{\EFourFigDir/fig_gain.pdf}
\caption{Speedup of the best tested configuration over the scenario's default
on the training grid. Tuning is worth most at small shapes and when the graph
is partitioned; it fades as the GPU becomes the bottleneck.}
\label{fig:e4-gain}
\end{figure*}

Fig.~\ref{fig:e4-gain} shows the same thing from the point of view of tuning.
The \emph{size} of the gain is strongly workload dependent, from 5.0$\times$
at $(1,128)$ with LayerNorm excluded down to a few percent at $(16,512)$. The
\emph{identity} of the best configuration is not, and only the latter is
something a selector can exploit.

\subsubsection{Selection on held-out shapes}

\begin{table}[t]
\caption{Held-out shapes: regret $R$ from~\eqref{eq:regret} and speedup over
default Torch-TensorRT. GELU and LayerNorm denote the scenarios in which that
operator runs in PyTorch.}
\label{tab:e4-heldout}
\centering
\footnotesize
\setlength{\tabcolsep}{2.6pt}
\begin{tabular}{@{}lrrrrrr@{}}
\toprule
 & \multicolumn{2}{c}{Unmodified} & \multicolumn{2}{c}{GELU}
 & \multicolumn{2}{c}{LayerNorm} \\
\cmidrule(lr){2-3}\cmidrule(lr){4-5}\cmidrule(l){6-7}
Method & $R$\,(\%) & speedup & $R$\,(\%) & speedup & $R$\,(\%) & speedup \\
\midrule
PyTorch eager        & 209.8 & 0.36$\times$ & 180.9 & 0.57$\times$ & 153.8 & 0.92$\times$ \\
Default TRT          &  11.0 & 1.00$\times$ &  61.3 & 1.00$\times$ & 132.7 & 1.00$\times$ \\
Fixed tuned          &   0.0 & 1.11$\times$ &   0.0 & 1.61$\times$ &   0.5 & 2.32$\times$ \\
Lookup sel.          &   0.0 & 1.11$\times$ &   3.2 & 1.56$\times$ &   1.4 & 2.29$\times$ \\
Regression sel.      &   1.5 & 1.09$\times$ &   0.5 & 1.60$\times$ &   0.0 & 2.33$\times$ \\
Best tested          &   0.0 & 1.11$\times$ &   0.0 & 1.61$\times$ &   0.0 & 2.33$\times$ \\
\bottomrule
\end{tabular}
\end{table}

\begin{figure}[t]
\centering
\includegraphics[width=\columnwidth]{\EFourFigDir/fig_regret.pdf}
\caption{Regret on the eight held-out shapes for the three methods that use
tuning data. The dashed line is the timing noise (0.7\%). No selector beats
the fixed tuned configuration by more than the noise.}
\label{fig:e4-regret}
\end{figure}

Table~\ref{tab:e4-heldout} and Fig.~\ref{fig:e4-regret} give the main result.
In all three scenarios the fixed tuned configuration is the CUDA Graphs
variant, and it is the best tested configuration on 23 of the 24
scenario--shape combinations. Its regret is 0.0\%, 0.0\% and 0.5\%. The
selectors range from 0.0\% to 3.2\%; relative to the fixed tuned configuration
they achieve between 0.97$\times$ and 1.01$\times$. The single case in which a
selector is ahead (regression, LayerNorm scenario, by 0.5\%) is below the
noise floor. Leave-one-out evaluation on the grid agrees: the fixed tuned
configuration has a regret of 0.0\%, 0.1\% and 0.9\%, the lookup selector
makes identical choices, and the regression selector scores 0.8\%, 1.0\% and
0.0\%.

Against default Torch-TensorRT, on the other hand, the gain is real and
sizeable: 1.11$\times$ on the unmodified model, 1.61$\times$ with GELU
excluded and 2.32$\times$ with LayerNorm excluded.

\subsubsection{Why there is nothing to select}

\begin{figure}[t]
\centering
\includegraphics[width=\columnwidth]{\EFourFigDir/fig_choice.pdf}
\caption{Fastest configuration at each of the 20 shapes in the two
partitioned scenarios. The only shape dependence is a switch to eager PyTorch
at the largest shapes. Circles mark held-out shapes where the lookup selector
chose differently from the best configuration, with the resulting loss.}
\label{fig:e4-choice}
\end{figure}

On the unmodified model, \cfg{trt\_cudagraphs} is the fastest configuration on
all 20 shapes. In the partitioned scenarios the CUDA Graphs variant is fastest
everywhere except at the top-right corner of the shape space
(Fig.~\ref{fig:e4-choice}): at $(16,512)$ in both scenarios and at $(16,448)$
with LayerNorm excluded, eager PyTorch is faster by 1\% to 11\%, because at
that size the per-call overhead no longer matters and the compiled module
pays for crossing between engines and PyTorch.

That one crossover is the entire workload dependence in our data, and the
selectors handle it poorly. Across the three scenarios the two selectors chose
eager for a held-out shape seven times, and it was the right call twice. The
lookup selector assigns $(12,384)$ to its nearest grid neighbor $(16,512)$ and
therefore to eager, at a cost of 23\% on that shape with GELU excluded and
12\% with LayerNorm excluded. The regression selector does better on the
partitioned scenarios but sends $(16,448)$ to eager on the unmodified model,
where the engine is 13\% faster. The underlying problem is that the boundary
is defined by a single training point at the edge of the grid, and that the
latency model is least accurate for exactly the configuration on the other
side of it: its mean absolute error on held-out shapes is 3\% to 4\% for the
GPU-bound configurations but 13\% for eager, whose flat-then-rising curve the
four terms of~\eqref{eq:model} do not fit well.

\subsubsection{Tuning cost}

\begin{figure*}[t]
\centering
\includegraphics[width=\textwidth]{\EFourFigDir/fig_budget.pdf}
\caption{Held-out regret when only $k$ of the twelve grid shapes are measured
(mean over subsets). Two shapes are enough to identify the fixed tuned
configuration. With one shape the outcome depends on which: if it is
$(16,512)$, every method selects eager PyTorch.}
\label{fig:e4-budget}
\end{figure*}

\begin{table}[t]
\caption{Tuning cost and break-even against default Torch-TensorRT, per
scenario}
\label{tab:e4-cost}
\centering
\footnotesize
\setlength{\tabcolsep}{4pt}
\begin{tabular}{@{}lrrr@{}}
\toprule
 & Unmodified & GELU & LayerNorm \\
\midrule
Saving per request (ms)        & 0.12 & 0.70 & 1.76 \\
\midrule
Full grid, 12 shapes (GPU-min) & 16.6 & 30.1 & 33.3 \\
\quad break-even (requests)    & $8.3\times10^6$ & $2.6\times10^6$ & $1.1\times10^6$ \\
\midrule
Two shapes (GPU-min)           & 2.7 & 5.1 & 5.5 \\
\quad break-even (requests)    & $1.4\times10^6$ & $4.3\times10^5$ & $1.9\times10^5$ \\
\bottomrule
\end{tabular}
\end{table}

The cost of a measurement is the engine build, the first call, the warm-up
and the timing. Measuring every configuration of a scenario on all twelve
grid shapes costs 17 to 33 GPU-minutes (Table~\ref{tab:e4-cost}), of which
engine builds are the bulk: a single-engine build takes 8 to 14\,s at the
default optimization level and 32\,s at level~5, a partitioned one about
30\,s, and the 13-engine \cfg{ln\_mbs1} about 57\,s.
Fig.~\ref{fig:e4-budget} shows that most of this effort is unnecessary. Two
measured shapes identify the fixed tuned configuration in every scenario, and
the ten further shapes buy nothing for it; the selectors need more data and
still end up level or behind. Dividing the cost by the mean saving per request
on the held-out shapes gives the number of requests after which tuning has
paid for itself: about 1.4 million on the unmodified model, where the saving
is 0.12\,ms, and 0.2 to 0.4 million in the partitioned scenarios.

% -----------------------------------------------------------------------------
\subsection{Discussion}
\label{sec:e4-discussion}

The proposal's hypothesis for this experiment was that a lookup table or a
regression model would recover most of the gap between a fixed configuration
and the best one. For this model the hypothesis fails for an uninteresting
reason: the gap is not there. We think the reason generalizes somewhat beyond
our setup, and it is worth stating.

A selector needs configurations whose latency curves \emph{cross}. In our
space, the configurations differ in how much per-call overhead they carry, and
overhead is a constant offset: a configuration with less of it is faster at
every shape, by an amount that shrinks as the shape grows but does not change
sign. CUDA Graphs removes the most overhead and therefore wins everywhere it
is applicable. The only crossing we observe involves eager PyTorch at the
largest shapes, and it is shallow (at most 11\%) and located at the edge of
the measured region. Per-workload selection becomes interesting when
configurations trade a fixed cost against a marginal one, for instance
precision modes with different kernel speeds, engines built with dynamic
versus static shape profiles, or models where a fallback block does real
computation instead of just interrupting the engine. None of these is in our
space.

For the project as a whole, three findings from this experiment matter more
than the selector:
\begin{enumerate}
  \item DistilBERT converts to a single TensorRT engine. Any partitioning
        effect studied on it has to be induced.
  \item On this model the cost of a partition boundary is launch overhead,
        and CUDA Graphs recovers most of it. Whether a graph is ``well
        partitioned'' matters far less with that option enabled.
  \item The shared benchmark's 20-iteration warm-up is too short for
        sub-millisecond engines on this hardware and overstates their latency
        by about 20\% in some repeats. Experiments that compare a single
        engine against PyTorch with it will understate the speedup.
\end{enumerate}

% -----------------------------------------------------------------------------
\subsection{Threats to Validity}
\label{sec:e4-threats}

\emph{The partitioned scenarios are artificial.} GELU and LayerNorm convert
without difficulty, and nobody would exclude them in practice. The scenarios
reproduce the structure of a fallback (execution leaves the engine and comes
back) but not necessarily its cost profile, since a genuinely unsupported
operator may be expensive in itself.

\emph{One model, one GPU, one precision.} Everything here is DistilBERT in
FP16 on one A100 node, with sequences of at most 512 tokens. All 340 pairs
were measured in a single job on a single machine, which removes
machine-to-machine variation from the comparison but also means we have not
observed it.

\emph{Static shapes.} Every engine is built for exactly one shape. A
deployment that serves many shapes from one engine with a dynamic profile
faces a different trade-off, which this experiment does not measure.

\emph{Drift in the CPU-bound baselines.} The held-out shapes were measured
in the last hour of the job, and the CPU-bound configurations look slower
there than on the grid. Relative to the grid-fitted model~\eqref{eq:model},
eager PyTorch is 15\% slower on the held-out shapes on average, the
LayerNorm-partitioned configurations 8\% to 11\%, and the GPU-bound
configurations 1\%. Some of this is model error, but neighboring shapes point
the same way (eager takes 3.25\,ms at $(8,256)$ on the grid and 4.06\,ms at
$(8,192)$ held out), so host-side performance probably degraded during the
job. Two consequences follow. The held-out speedups over eager and over the
partitioned defaults are somewhat inflated; the leave-one-out figures on the
grid (2.22$\times$ over the default for LayerNorm, against 2.32$\times$ held
out) are the more conservative ones. And the position of the eager crossover
on the held-out shapes is uncertain: if eager were in fact 15\% faster at
$(16,448)$, it would be the best choice there in all three scenarios, the
fixed tuned configuration would lose 2\% to 20\% on that one shape, and a
selector choosing eager there would be right. This would move the
held-out regret of the fixed tuned configuration to at most 2.3\%. It would not
change the picture at the other seven shapes or on the grid.

\emph{Build times share a cache.} TensorRT reuses layer timings within a
process, so the first build of the job took 40\,s and later builds of the same
graph 6 to 13\,s. The tuning costs in Table~\ref{tab:e4-cost} are therefore
optimistic for a cold start.

\emph{Memory.} The recorded peak memory counts PyTorch's allocator only and
excludes memory held by TensorRT engines, so we do not report a memory
comparison.

% -----------------------------------------------------------------------------
\subsection{Summary}
\label{sec:e4-summary}

We built a resumable sweep harness, two workload-aware selectors and an
evaluation on held-out shapes, and measured 17 configurations on 20 shapes.
For DistilBERT on an A100, tuning the configuration once is worth
1.11$\times$ to 2.32$\times$ over default Torch-TensorRT, two measured shapes
are enough to do it, and selecting the configuration per workload adds nothing
beyond that. The result rests on a measurement protocol that had to be
corrected first; with the original warm-up the same experiment would have
reported a workload-dependent optimum that does not exist.
 from the team's main file.
% Needs in the preamble: graphicx, booktabs, amsmath, cite, url.
% Figures are read from \EFourFigDir (default: figures/). Redefine it before the
% \input if the figures live elsewhere, e.g. \newcommand{\EFourFigDir}{e4/figures}.
% References are in e4_refs.bib.
% =============================================================================
\providecommand{\EFourFigDir}{figures}
\providecommand{\cfg}[1]{\texttt{#1}}

\section{Workload-Aware Configuration Selection}
\label{sec:e4}

The previous experiments establish that Torch-TensorRT latency depends on how
the graph is partitioned and on the shape of the request. This experiment asks
the question that follows from both: if the best compiler configuration moves
with the workload, can a cheap model learn where it moves, so that a deployment
picks its configuration from the batch size and sequence length instead of
committing to one setting for everything? Rice's algorithm selection
problem~\cite{rice1976algorithm} gives the framing, and cost-model-driven
autotuners such as AutoTVM~\cite{chen2018autotvm} show that the idea pays off
when the configuration space is large and the optimum is workload dependent.

Our answer for DistilBERT on an A100 is negative, and we report it as such. A
single tuned configuration is within 0.5\% of the best tested configuration on
every held-out shape. The two selectors we built reproduce that configuration
almost everywhere and, where they deviate, they mostly lose. What does matter
is tuning at all: the tuned configuration is 1.11$\times$ faster than default
Torch-TensorRT on the unmodified model and 1.6$\times$ to 2.3$\times$ faster
once part of the graph falls back to PyTorch. Getting to a result we trust took
three iterations, and two of them changed the experiment more than the
selector did. We describe them in the order they happened, because the
intermediate findings (a model with nothing to partition, a cluster that
cancels compile-heavy jobs, and a 20\% timing artifact) are useful to the rest
of the project on their own.

% -----------------------------------------------------------------------------
\subsection{Problem Formulation}
\label{sec:e4-problem}

A workload is a request shape $w=(b,s)$, with batch size $b$ and sequence
length $s$ in tokens. A configuration $c \in \mathcal{C}$ is one way of running
the model: PyTorch eager, or Torch-TensorRT with a particular set of compiler
and runtime settings. $L(c,w)$ is the warmed-up median latency of $c$ on $w$. A
selector is a function $\pi:\,w \mapsto c$ fitted on measurements from a set of
training shapes $W_\mathrm{train}$ and judged on shapes it has not seen. For a
test set $W$ we report the \emph{regret} of $\pi$ against the best tested
configuration,
\begin{equation}
  R(\pi) \;=\;
  \Bigl(\,\prod_{w\in W}
     \frac{L(\pi(w),w)}{\min_{c\in\mathcal{C}} L(c,w)}\Bigr)^{1/|W|} - 1 ,
  \label{eq:regret}
\end{equation}
that is, the geometric-mean slowdown relative to an oracle that tries every
configuration on every shape. A geometric mean keeps a 5\% loss at batch~1 and
a 5\% loss at batch~16 on equal footing, although the second is ten times
larger in milliseconds.

We compare six methods, following the evaluation plan of the proposal:
\begin{itemize}
  \item \textbf{PyTorch eager}, the uncompiled FP16 model;
  \item \textbf{Default Torch-TensorRT}, the compiler's default settings;
  \item \textbf{One fixed tuned configuration}, the configuration with the
        lowest geometric-mean slowdown over $W_\mathrm{train}$, used for every
        shape;
  \item \textbf{Lookup selector}, which returns the best configuration of the
        nearest training shape, with distance measured in
        $(\log_2 b,\log_2 s)$;
  \item \textbf{Regression selector}, which fits one latency model per
        configuration and returns the configuration with the lowest
        prediction (Section~\ref{sec:e4-selectors});
  \item \textbf{Best tested configuration}, the oracle in~\eqref{eq:regret}.
\end{itemize}
The fixed tuned configuration is the baseline that decides the experiment. A
selector is only worth deploying if it beats the best single configuration,
not merely the default one.

% -----------------------------------------------------------------------------
\subsection{Experimental Setup}
\label{sec:e4-setup}

All measurements come from one NVIDIA A100-SXM4-40GB node of the NYU HPC cloud
bursting cluster, inside a Singularity container with PyTorch~2.14.1
(CUDA~13.0)~\cite{paszke2019pytorch}, Torch-TensorRT~2.14.0~\cite{torchtrt} and
TensorRT~11.1.0~\cite{tensorrt}. The model is
\cfg{distilbert-base-uncased}~\cite{sanh2019distilbert} in FP16, wrapped so
that it takes token ids and an attention mask and returns the last hidden
state. Inputs are random token ids from a fixed seed with a full attention
mask, and every configuration receives the same tensors. Torch-TensorRT uses
its Dynamo front end, and each engine is built for one static shape, so a
(configuration, shape) pair is one compilation.

The training grid is $b\in\{1,4,8,16\}\times s\in\{128,256,512\}$, twelve
shapes. The proposal also lists 1024 tokens, but DistilBERT has 512 position
embeddings, so longer inputs cannot be run without changing the model. Eight
further shapes lie off the grid and are never used for fitting:
$(1,384)$, $(2,192)$, $(2,384)$, $(4,320)$, $(8,192)$, $(12,192)$, $(12,384)$
and $(16,448)$. None of their sequence lengths occurs in the grid, and four
of them also use a batch size (2 or 12) that does not.

Latency is measured with the project's shared benchmark routine: 20 warm-up
calls followed by 100 calls timed with CUDA events, from which we take the
median and the 95th percentile. Around it, our sweep harness records the
compile time, the time of the first call, the number of TensorRT and PyTorch
blocks in the compiled module, and the agreement of the output with eager
PyTorch (cosine similarity and largest absolute difference). The harness
appends one row per timing repeat to a CSV file and skips pairs that are
already present, so an interrupted job resumes where it stopped. That property
turned out to be necessary (Section~\ref{sec:e4-iter2}).

\begin{table}[t]
\caption{Chronology of cluster runs for this experiment (all on 12 Nov.\ 2026)}
\label{tab:e4-runs}
\centering
\footnotesize
\begin{tabular}{@{}llp{4.55cm}@{}}
\toprule
Job & Purpose & Outcome \\
\midrule
29499 & Environment & Cancelled: \cfg{--fakeroot} unpacked the 8.5\,GB image to
        network storage \\
29507 & Environment & Held by the scheduler; resubmitted \\
29509 & Environment & Packages installed (plain writable overlay) \\
29537 & Partition check & One engine, no fallback; LayerNorm exclusion had no
        effect (wrong operator name) \\
29550 & Partition check & Corrected; Table~\ref{tab:e4-partition} \\
29551 & Sweep 1 & GPU node never booted; cancelled after 12\,min \\
29586 & Sweep 1 & 83 of 340 runs, then cancelled by the cluster's
        low-utilization watchdog at 30\,min \\
29727--8 & Sweep 2 & Two more node boot failures (13\,min each) \\
29729 & Sweep 2 & 340 of 340 runs in 2\,h\,31\,min; all results below \\
\bottomrule
\end{tabular}
\end{table}

Table~\ref{tab:e4-runs} lists every cluster job behind this section. Only the
last one contributes numbers to the results; the others are the reason the
last one looks the way it does.

% -----------------------------------------------------------------------------
\subsection{Iteration 1: Is There Anything to Tune?}
\label{sec:e4-iter1}

Our plan assumed that the partitioning controls of Torch-TensorRT
(\cfg{min\_block\_size}, the choice between the fast and the global
partitioner, and operator exclusion) would produce visibly different graphs
for DistilBERT, and that the selector would arbitrate between them. Before
spending GPU hours on engines we checked that assumption with a dry-run
compile of every candidate at $(4,128)$, which reports the partition without
building anything.

The assumption was wrong. After lowering, DistilBERT consists of 127
operators (36 linear layers, 24 reshapes, 24 permutes, 13 additions, 13 layer
normalizations, six attention and six GELU operators, and a handful of
embedding and mask operations), and TensorRT supports all of them. The default
compile is therefore a single engine with no PyTorch fallback, and
\cfg{min\_block\_size} and the partitioner have nothing to act on: we verified
that \cfg{min\_block\_size}$\,=1$ and the global partitioner return the
identical one-engine graph, and removed both from the sweep.

That left two options: switch to a model with unsupported operators, which
would have cut this experiment off from the shared harness and from the other
three experiments, or keep the model and create the partitioning problem
deliberately. We chose the second and say so plainly, because it limits what
the results mean (Section~\ref{sec:e4-threats}). Torch-TensorRT lets the user
force named operators to run in PyTorch. We define three \emph{scenarios}:
\begin{itemize}
  \item \textbf{Unmodified model}: nothing excluded, one engine.
  \item \textbf{GELU in PyTorch}: the six GELU operators are excluded, which
        splits the encoder between its layers.
  \item \textbf{LayerNorm in PyTorch}: the thirteen layer normalizations are
        excluded, which cuts every layer twice.
\end{itemize}
Each scenario is evaluated separately, with its own default, its own tuned
configuration and its own selectors, as if it were a different model.

\begin{table}[t]
\caption{Partition structure of each configuration (dry-run compile)}
\label{tab:e4-partition}
\centering
\footnotesize
\setlength{\tabcolsep}{3.5pt}
\begin{tabular}{@{}lccl@{}}
\toprule
 & \multicolumn{2}{c}{Blocks} & \\
\cmidrule(lr){2-3}
Configuration & TRT & PyT. & Operators per TRT block \\
\midrule
\cfg{trt\_default}, \cfg{trt\_opt*} & 1 & 0 & 127 \\
\midrule
\cfg{gelu\_mbs5}, \cfg{gelu\_global} & 6 & 6 & 23, then $5\times19$ \\
\cfg{gelu\_mbs1} & 7 & 6 & 23, $5\times19$, 3 \\
\cfg{gelu\_mbs22} & 1 & 1 & 23 (104 ops in PyTorch) \\
\midrule
\cfg{ln\_mbs5} & 7 & 7 & 6, then $6\times14$ \\
\cfg{ln\_mbs1} & 13 & 13 & 6, then $6\times(14,\,4)$ \\
\cfg{ln\_mbs10} & 6 & 7 & $6\times14$ \\
\cfg{ln\_mbs20} & 0 & 1 & none (127 ops in PyTorch) \\
\cfg{ln\_global} & 6 & 7 & 18, then $5\times14$ \\
\bottomrule
\end{tabular}
\end{table}

Table~\ref{tab:e4-partition} shows the resulting partitions. The dry run also
pruned the configuration list a second time. With GELU excluded, the TensorRT
blocks hold 19 to 23 operators, so \cfg{min\_block\_size} values of 5, 10 and
20 all give the same graph; we kept 5 and replaced the other two by 22, the
smallest value that changes anything. Two mistakes surfaced in the same pass.
Our first LayerNorm scenario excluded \cfg{aten.layer\_norm}, but the operator
that exists after lowering is \cfg{aten.native\_layer\_norm}; the exclusion
silently matched nothing and the ``broken'' model was still one engine. And
our block counter reported zero PyTorch blocks for the global partitioner,
because that partitioner leaves fallback operators inline in the top-level
graph instead of wrapping them in submodules. Both were fixed before any
latency was measured. Without the dry run, a third of the sweep would have
been mislabeled.

\begin{table}[t]
\caption{Configurations in the sweep. Default settings are
\cfg{min\_block\_size}$\,=5$, fast partitioner, optimization level~3.}
\label{tab:e4-configs}
\centering
\footnotesize
\setlength{\tabcolsep}{3.5pt}
\begin{tabular}{@{}lll@{}}
\toprule
Scenario & Configuration & Change from scenario default \\
\midrule
all & \cfg{eager} & PyTorch, no compilation \\
\midrule
Unmodified & \cfg{trt\_default} & (default) \\
 & \cfg{trt\_opt1} & optimization level 1 \\
 & \cfg{trt\_opt5} & optimization level 5 \\
 & \cfg{trt\_cudagraphs} & CUDA Graphs at run time \\
 & \cfg{trt\_default\_rebuild} & none (second build) \\
\midrule
GELU in & \cfg{gelu\_mbs5} & (default) \\
PyTorch & \cfg{gelu\_mbs1} & \cfg{min\_block\_size} 1 \\
 & \cfg{gelu\_mbs22} & \cfg{min\_block\_size} 22 \\
 & \cfg{gelu\_global} & global partitioner \\
 & \cfg{gelu\_cudagraphs} & CUDA Graphs at run time \\
\midrule
LayerNorm & \cfg{ln\_mbs5} & (default) \\
in PyTorch & \cfg{ln\_mbs1} & \cfg{min\_block\_size} 1 \\
 & \cfg{ln\_mbs10}, \cfg{ln\_mbs20} & \cfg{min\_block\_size} 10, 20 \\
 & \cfg{ln\_global} & global partitioner \\
 & \cfg{ln\_cudagraphs} & CUDA Graphs at run time \\
\bottomrule
\end{tabular}
\end{table}

The final space has 17 configurations (Table~\ref{tab:e4-configs}). Because
partitioning offers nothing on the unmodified model, that scenario varies the
builder's optimization level instead. All three scenarios include whole-graph
CUDA Graphs~\cite{gray2019cudagraphs}, a runtime option that records the
module's stream of GPU kernel launches once and replays it as a single
submission. One more configuration, \cfg{trt\_default\_rebuild}, is the
default built a second time. It belongs to no scenario and exists to measure
how much two builds of the same configuration differ.

% -----------------------------------------------------------------------------
\subsection{Iteration 2: A First Sweep That Could Not Be Used}
\label{sec:e4-iter2}

The first sweep timed each compiled module three times with the shared
protocol. It was cancelled after 30 minutes with 83 of 340 runs done, and when
we looked at those 83 runs they were not usable either. The two problems are
unrelated.

\subsubsection{Compile-heavy jobs are cancelled}
The accounting record showed the job cancelled by the administrator account
with no message. The cause is a watchdog on the cluster that cancels a GPU job
once it is about 30 minutes old if its mean GPU utilization is 5\% or less. A
configuration sweep is exactly such a job: more than 85\% of its wall time is
TensorRT building engines, which is mostly CPU work, and timing a sub-millisecond engine
for 100 iterations keeps the GPU busy for a tenth of a second. We did not add
artificial load to stay above the threshold. The harness was already
resumable, so the second sweep was submitted as a chain of jobs, each starting
when the previous one ends and continuing from the same file.

\subsubsection{A 20\% timing artifact}
\label{sec:e4-warmup}

\begin{figure}[t]
\centering
\includegraphics[width=\columnwidth]{\EFourFigDir/fig_warmup.pdf}
\caption{Each timing repeat relative to the fastest repeat of the same
compiled module, for the five shapes both sweeps cover. (a)~With the shared
20-iteration warm-up, GPU-bound configurations start about 20\% slow and drop
to their final level one or two repeats later. (b)~With two seconds of
sustained load before timing, all repeats agree.}
\label{fig:e4-warmup}
\end{figure}

For 30 of the 79 successful runs, the three repeats of the \emph{same engine}
disagreed by more than 8\%. The disagreement has a clear shape
(Fig.~\ref{fig:e4-warmup}a): the first one or two repeats are 20\% to 25\%
slower, then the latency drops and stays down. It never goes the other way. It
affects only configurations whose latency is set by the GPU, namely the single
engine and everything run under CUDA Graphs; eager PyTorch and the
configurations with PyTorch fallback blocks, whose latency is set by the CPU
launching kernels, are stable to within 4\%.

Our reading is that the GPU needs a few hundred milliseconds of sustained
load to reach full speed after sitting idle through a compile, and that the
20-iteration warm-up is over long before that when an iteration takes
0.6\,ms. We did not confirm the mechanism (clock or power state) directly, so
we treat it as an observation about this node. The effect on a selector study
is severe either way. The differences we are trying to resolve on the
unmodified model are of the order of 10\%, and with three repeats the median
lands on the slow or the fast level essentially at random. Under these
conditions the selector would have been fitted to noise, and it would have
looked as if the best configuration changed from shape to shape. That
steady-state measurement is easy to get wrong without doing anything obviously
wrong is well documented~\cite{georges2007statistically,mytkowicz2009wrong};
this is one more instance.

% -----------------------------------------------------------------------------
\subsection{Iteration 3: The Corrected Sweep}
\label{sec:e4-iter3}

The second sweep differs from the first in three ways. Before timing, the
harness now keeps the GPU busy with the workload itself for two seconds; the
shared benchmark routine then runs unchanged, including its own 20 warm-up
calls. Each module is timed five times instead of three, and the latency we
use is the median of the five medians. And every row records the node and job
it came from, in case a chained sweep ends up spread over several machines.
The first sweep's data was discarded.

After two more failed node boots, the third job of the chain ran all 340 pairs
in one piece, in 2\,h\,31\,min on a single node. The warm-up had an unplanned
side effect: it lifted the job's GPU utilization above the watchdog threshold,
so the chain was never needed.

The data is clean. Of the 340 pairs, 320 succeeded. The 20 failures are all
\cfg{ln\_global}, which fails to build on every shape with an internal error
in Torch-TensorRT 2.14 (a missing \cfg{\_frozen\_param0} attribute on the
partitioned graph); it is excluded from selection. Across the five repeats the
median spread is 0.7\%, the 90th percentile 2.4\% and the maximum 5.8\%
(Fig.~\ref{fig:e4-warmup}b). The default configuration and its rebuild differ
by 0.3\% in the median. We therefore treat differences below about 1\% as
ties. Every compiled module agrees with eager PyTorch to a cosine similarity
of at least 0.99999, with a largest absolute difference of 0.02 in FP16.

% -----------------------------------------------------------------------------
\subsection{Selectors}
\label{sec:e4-selectors}

Both selectors are deliberately small. With twelve training shapes there is no
room for anything else.

The \emph{lookup selector} stores, for each training shape, the configuration
that was fastest there, and answers a query with the entry of the nearest
training shape in $(\log_2 b, \log_2 s)$. It has no parameters and cannot
extrapolate.

The \emph{regression selector} fits, for each configuration $c$, the latency
model
\begin{equation}
  \hat{L}_c(b,s) \;=\; \theta_{c,0} + \theta_{c,1}\,b + \theta_{c,2}\,bs
                      + \theta_{c,3}\,bs^2 ,
  \qquad \theta_c \ge 0 ,
  \label{eq:model}
\end{equation}
whose terms stand for a fixed per-call overhead, a per-sequence cost, a
per-token cost and the quadratic cost of self-attention. The coefficients are
obtained by non-negative least squares~\cite{lawson1995solving} on the
\emph{relative} error, $\sum_{w}(\hat{L}_c(w)/L(c,w)-1)^2$, so that small and
large shapes carry equal weight. Non-negativity keeps the model physically
meaningful and prevents it from predicting negative latencies outside the
grid. The selector returns $\arg\min_c \hat{L}_c(b,s)$.

We evaluate in two ways. \emph{Held-out}: fit on the twelve grid shapes, test
on the eight off-grid shapes. \emph{Leave-one-out}: for each grid shape, fit
on the other eleven and test on the one left out. To relate quality to tuning
effort, we also refit every method on subsets of $k=1,\dots,12$ grid shapes
(all subsets where there are at most 40, otherwise 40 random ones) and test on
the held-out shapes.

% -----------------------------------------------------------------------------
\subsection{Results}
\label{sec:e4-results}

\subsubsection{The latency landscape}

\begin{figure*}[t]
\centering
\includegraphics[width=\textwidth]{\EFourFigDir/fig_latency.pdf}
\caption{Median latency against batch size at 512 tokens. Eager PyTorch is
flat until the largest batch, which shows that it is limited by per-call
overhead and not by GPU work. A partitioned model inherits part of that
overhead (circles, squares); CUDA Graphs removes it (triangles). All curves meet at
batch 16, where GPU work dominates.}
\label{fig:e4-latency}
\end{figure*}

Fig.~\ref{fig:e4-latency} shows the structure that everything else follows
from. Eager PyTorch takes between 3.2 and 4.1\,ms on every shape below
$(16,448)$, almost independent of size: at these sizes an FP16 DistilBERT
forward pass on an A100 is bounded by the time it takes Python and the
framework to issue a long sequence of small kernels, not by the kernels
themselves. The single
TensorRT engine removes that overhead and is 2.68$\times$ faster than eager in
the geometric mean over all 20 shapes, from 6.6$\times$ at $(1,128)$ down to
1.07$\times$ at $(16,512)$, where the GPU is finally the bottleneck.

A partitioned model sits in between, and where it sits depends on how many
times execution returns to PyTorch. With GELU excluded (six fallback blocks)
the default configuration is 1.63$\times$ slower than the single engine; with
LayerNorm excluded (seven) it is 2.49$\times$ slower and only 7\% faster than
not compiling at all. Table~\ref{tab:e4-knobs} summarizes what each control is
worth.

\begin{table}[t]
\caption{Effect of each control on latency (geometric mean over all 20 shapes,
range over shapes in parentheses)}
\label{tab:e4-knobs}
\centering
\footnotesize
\begin{tabular}{@{}p{4.35cm}l@{}}
\toprule
Change & Latency \\
\midrule
\multicolumn{2}{@{}l}{\emph{Unmodified model, relative to default}} \\
Optimization level 5 & 1.00$\times$ (0.97--1.02) \\
Optimization level 1 & 1.06$\times$ slower (1.03--1.08) \\
CUDA Graphs & 1.11$\times$ faster (1.02--1.20) \\
\midrule
\multicolumn{2}{@{}l}{\emph{GELU in PyTorch, relative to its default}} \\
\cfg{min\_block\_size} 1 (7 engines) & 1.01$\times$ faster (0.99--1.04) \\
Global partitioner & 1.02$\times$ slower (0.98--1.05) \\
\cfg{min\_block\_size} 22 (1 engine) & 1.74$\times$ slower (1.07--2.25) \\
CUDA Graphs & 1.63$\times$ faster (1.06--3.27) \\
\midrule
\multicolumn{2}{@{}l}{\emph{LayerNorm in PyTorch, relative to its default}} \\
\cfg{min\_block\_size} 1 (13 engines) & 1.05$\times$ faster (0.99--1.11) \\
\cfg{min\_block\_size} 10 (6 engines) & 1.03$\times$ slower (1.00--1.05) \\
\cfg{min\_block\_size} 20 (no engine) & 1.28$\times$ slower (1.04--1.47) \\
CUDA Graphs & 2.26$\times$ faster (1.08--5.02) \\
\bottomrule
\end{tabular}
\end{table}

Two things stand out. First, the partitioning controls barely matter as long
as the engines survive: \cfg{min\_block\_size} and the partitioner move
latency by 5\% or less, and their only large effect is the destructive one,
when a high threshold discards engines and pushes the model back towards
eager. Second, one runtime option dwarfs all of them. CUDA Graphs makes the
GELU-partitioned model 1.63$\times$ faster and the LayerNorm-partitioned one
2.26$\times$ faster, to within 11\% and 23\% of the single engine under the
same option. This is consistent with the picture above: what a fallback block
costs is launch overhead, and replaying a recorded launch sequence removes it
no matter how many blocks there are.

\begin{figure*}[t]
\centering
\includegraphics[width=\textwidth]{\EFourFigDir/fig_gain.pdf}
\caption{Speedup of the best tested configuration over the scenario's default
on the training grid. Tuning is worth most at small shapes and when the graph
is partitioned; it fades as the GPU becomes the bottleneck.}
\label{fig:e4-gain}
\end{figure*}

Fig.~\ref{fig:e4-gain} shows the same thing from the point of view of tuning.
The \emph{size} of the gain is strongly workload dependent, from 5.0$\times$
at $(1,128)$ with LayerNorm excluded down to a few percent at $(16,512)$. The
\emph{identity} of the best configuration is not, and only the latter is
something a selector can exploit.

\subsubsection{Selection on held-out shapes}

\begin{table}[t]
\caption{Held-out shapes: regret $R$ from~\eqref{eq:regret} and speedup over
default Torch-TensorRT. GELU and LayerNorm denote the scenarios in which that
operator runs in PyTorch.}
\label{tab:e4-heldout}
\centering
\footnotesize
\setlength{\tabcolsep}{2.6pt}
\begin{tabular}{@{}lrrrrrr@{}}
\toprule
 & \multicolumn{2}{c}{Unmodified} & \multicolumn{2}{c}{GELU}
 & \multicolumn{2}{c}{LayerNorm} \\
\cmidrule(lr){2-3}\cmidrule(lr){4-5}\cmidrule(l){6-7}
Method & $R$\,(\%) & speedup & $R$\,(\%) & speedup & $R$\,(\%) & speedup \\
\midrule
PyTorch eager        & 209.8 & 0.36$\times$ & 180.9 & 0.57$\times$ & 153.8 & 0.92$\times$ \\
Default TRT          &  11.0 & 1.00$\times$ &  61.3 & 1.00$\times$ & 132.7 & 1.00$\times$ \\
Fixed tuned          &   0.0 & 1.11$\times$ &   0.0 & 1.61$\times$ &   0.5 & 2.32$\times$ \\
Lookup sel.          &   0.0 & 1.11$\times$ &   3.2 & 1.56$\times$ &   1.4 & 2.29$\times$ \\
Regression sel.      &   1.5 & 1.09$\times$ &   0.5 & 1.60$\times$ &   0.0 & 2.33$\times$ \\
Best tested          &   0.0 & 1.11$\times$ &   0.0 & 1.61$\times$ &   0.0 & 2.33$\times$ \\
\bottomrule
\end{tabular}
\end{table}

\begin{figure}[t]
\centering
\includegraphics[width=\columnwidth]{\EFourFigDir/fig_regret.pdf}
\caption{Regret on the eight held-out shapes for the three methods that use
tuning data. The dashed line is the timing noise (0.7\%). No selector beats
the fixed tuned configuration by more than the noise.}
\label{fig:e4-regret}
\end{figure}

Table~\ref{tab:e4-heldout} and Fig.~\ref{fig:e4-regret} give the main result.
In all three scenarios the fixed tuned configuration is the CUDA Graphs
variant, and it is the best tested configuration on 23 of the 24
scenario--shape combinations. Its regret is 0.0\%, 0.0\% and 0.5\%. The
selectors range from 0.0\% to 3.2\%; relative to the fixed tuned configuration
they achieve between 0.97$\times$ and 1.01$\times$. The single case in which a
selector is ahead (regression, LayerNorm scenario, by 0.5\%) is below the
noise floor. Leave-one-out evaluation on the grid agrees: the fixed tuned
configuration has a regret of 0.0\%, 0.1\% and 0.9\%, the lookup selector
makes identical choices, and the regression selector scores 0.8\%, 1.0\% and
0.0\%.

Against default Torch-TensorRT, on the other hand, the gain is real and
sizeable: 1.11$\times$ on the unmodified model, 1.61$\times$ with GELU
excluded and 2.32$\times$ with LayerNorm excluded.

\subsubsection{Why there is nothing to select}

\begin{figure}[t]
\centering
\includegraphics[width=\columnwidth]{\EFourFigDir/fig_choice.pdf}
\caption{Fastest configuration at each of the 20 shapes in the two
partitioned scenarios. The only shape dependence is a switch to eager PyTorch
at the largest shapes. Circles mark held-out shapes where the lookup selector
chose differently from the best configuration, with the resulting loss.}
\label{fig:e4-choice}
\end{figure}

On the unmodified model, \cfg{trt\_cudagraphs} is the fastest configuration on
all 20 shapes. In the partitioned scenarios the CUDA Graphs variant is fastest
everywhere except at the top-right corner of the shape space
(Fig.~\ref{fig:e4-choice}): at $(16,512)$ in both scenarios and at $(16,448)$
with LayerNorm excluded, eager PyTorch is faster by 1\% to 11\%, because at
that size the per-call overhead no longer matters and the compiled module
pays for crossing between engines and PyTorch.

That one crossover is the entire workload dependence in our data, and the
selectors handle it poorly. Across the three scenarios the two selectors chose
eager for a held-out shape seven times, and it was the right call twice. The
lookup selector assigns $(12,384)$ to its nearest grid neighbor $(16,512)$ and
therefore to eager, at a cost of 23\% on that shape with GELU excluded and
12\% with LayerNorm excluded. The regression selector does better on the
partitioned scenarios but sends $(16,448)$ to eager on the unmodified model,
where the engine is 13\% faster. The underlying problem is that the boundary
is defined by a single training point at the edge of the grid, and that the
latency model is least accurate for exactly the configuration on the other
side of it: its mean absolute error on held-out shapes is 3\% to 4\% for the
GPU-bound configurations but 13\% for eager, whose flat-then-rising curve the
four terms of~\eqref{eq:model} do not fit well.

\subsubsection{Tuning cost}

\begin{figure*}[t]
\centering
\includegraphics[width=\textwidth]{\EFourFigDir/fig_budget.pdf}
\caption{Held-out regret when only $k$ of the twelve grid shapes are measured
(mean over subsets). Two shapes are enough to identify the fixed tuned
configuration. With one shape the outcome depends on which: if it is
$(16,512)$, every method selects eager PyTorch.}
\label{fig:e4-budget}
\end{figure*}

\begin{table}[t]
\caption{Tuning cost and break-even against default Torch-TensorRT, per
scenario}
\label{tab:e4-cost}
\centering
\footnotesize
\setlength{\tabcolsep}{4pt}
\begin{tabular}{@{}lrrr@{}}
\toprule
 & Unmodified & GELU & LayerNorm \\
\midrule
Saving per request (ms)        & 0.12 & 0.70 & 1.76 \\
\midrule
Full grid, 12 shapes (GPU-min) & 16.6 & 30.1 & 33.3 \\
\quad break-even (requests)    & $8.3\times10^6$ & $2.6\times10^6$ & $1.1\times10^6$ \\
\midrule
Two shapes (GPU-min)           & 2.7 & 5.1 & 5.5 \\
\quad break-even (requests)    & $1.4\times10^6$ & $4.3\times10^5$ & $1.9\times10^5$ \\
\bottomrule
\end{tabular}
\end{table}

The cost of a measurement is the engine build, the first call, the warm-up
and the timing. Measuring every configuration of a scenario on all twelve
grid shapes costs 17 to 33 GPU-minutes (Table~\ref{tab:e4-cost}), of which
engine builds are the bulk: a single-engine build takes 8 to 14\,s at the
default optimization level and 32\,s at level~5, a partitioned one about
30\,s, and the 13-engine \cfg{ln\_mbs1} about 57\,s.
Fig.~\ref{fig:e4-budget} shows that most of this effort is unnecessary. Two
measured shapes identify the fixed tuned configuration in every scenario, and
the ten further shapes buy nothing for it; the selectors need more data and
still end up level or behind. Dividing the cost by the mean saving per request
on the held-out shapes gives the number of requests after which tuning has
paid for itself: about 1.4 million on the unmodified model, where the saving
is 0.12\,ms, and 0.2 to 0.4 million in the partitioned scenarios.

% -----------------------------------------------------------------------------
\subsection{Discussion}
\label{sec:e4-discussion}

The proposal's hypothesis for this experiment was that a lookup table or a
regression model would recover most of the gap between a fixed configuration
and the best one. For this model the hypothesis fails for an uninteresting
reason: the gap is not there. We think the reason generalizes somewhat beyond
our setup, and it is worth stating.

A selector needs configurations whose latency curves \emph{cross}. In our
space, the configurations differ in how much per-call overhead they carry, and
overhead is a constant offset: a configuration with less of it is faster at
every shape, by an amount that shrinks as the shape grows but does not change
sign. CUDA Graphs removes the most overhead and therefore wins everywhere it
is applicable. The only crossing we observe involves eager PyTorch at the
largest shapes, and it is shallow (at most 11\%) and located at the edge of
the measured region. Per-workload selection becomes interesting when
configurations trade a fixed cost against a marginal one, for instance
precision modes with different kernel speeds, engines built with dynamic
versus static shape profiles, or models where a fallback block does real
computation instead of just interrupting the engine. None of these is in our
space.

For the project as a whole, three findings from this experiment matter more
than the selector:
\begin{enumerate}
  \item DistilBERT converts to a single TensorRT engine. Any partitioning
        effect studied on it has to be induced.
  \item On this model the cost of a partition boundary is launch overhead,
        and CUDA Graphs recovers most of it. Whether a graph is ``well
        partitioned'' matters far less with that option enabled.
  \item The shared benchmark's 20-iteration warm-up is too short for
        sub-millisecond engines on this hardware and overstates their latency
        by about 20\% in some repeats. Experiments that compare a single
        engine against PyTorch with it will understate the speedup.
\end{enumerate}

% -----------------------------------------------------------------------------
\subsection{Threats to Validity}
\label{sec:e4-threats}

\emph{The partitioned scenarios are artificial.} GELU and LayerNorm convert
without difficulty, and nobody would exclude them in practice. The scenarios
reproduce the structure of a fallback (execution leaves the engine and comes
back) but not necessarily its cost profile, since a genuinely unsupported
operator may be expensive in itself.

\emph{One model, one GPU, one precision.} Everything here is DistilBERT in
FP16 on one A100 node, with sequences of at most 512 tokens. All 340 pairs
were measured in a single job on a single machine, which removes
machine-to-machine variation from the comparison but also means we have not
observed it.

\emph{Static shapes.} Every engine is built for exactly one shape. A
deployment that serves many shapes from one engine with a dynamic profile
faces a different trade-off, which this experiment does not measure.

\emph{Drift in the CPU-bound baselines.} The held-out shapes were measured
in the last hour of the job, and the CPU-bound configurations look slower
there than on the grid. Relative to the grid-fitted model~\eqref{eq:model},
eager PyTorch is 15\% slower on the held-out shapes on average, the
LayerNorm-partitioned configurations 8\% to 11\%, and the GPU-bound
configurations 1\%. Some of this is model error, but neighboring shapes point
the same way (eager takes 3.25\,ms at $(8,256)$ on the grid and 4.06\,ms at
$(8,192)$ held out), so host-side performance probably degraded during the
job. Two consequences follow. The held-out speedups over eager and over the
partitioned defaults are somewhat inflated; the leave-one-out figures on the
grid (2.22$\times$ over the default for LayerNorm, against 2.32$\times$ held
out) are the more conservative ones. And the position of the eager crossover
on the held-out shapes is uncertain: if eager were in fact 15\% faster at
$(16,448)$, it would be the best choice there in all three scenarios, the
fixed tuned configuration would lose 2\% to 20\% on that one shape, and a
selector choosing eager there would be right. This would move the
held-out regret of the fixed tuned configuration to at most 2.3\%. It would not
change the picture at the other seven shapes or on the grid.

\emph{Build times share a cache.} TensorRT reuses layer timings within a
process, so the first build of the job took 40\,s and later builds of the same
graph 6 to 13\,s. The tuning costs in Table~\ref{tab:e4-cost} are therefore
optimistic for a cold start.

\emph{Memory.} The recorded peak memory counts PyTorch's allocator only and
excludes memory held by TensorRT engines, so we do not report a memory
comparison.

% -----------------------------------------------------------------------------
\subsection{Summary}
\label{sec:e4-summary}

We built a resumable sweep harness, two workload-aware selectors and an
evaluation on held-out shapes, and measured 17 configurations on 20 shapes.
For DistilBERT on an A100, tuning the configuration once is worth
1.11$\times$ to 2.32$\times$ over default Torch-TensorRT, two measured shapes
are enough to do it, and selecting the configuration per workload adds nothing
beyond that. The result rests on a measurement protocol that had to be
corrected first; with the original warm-up the same experiment would have
reported a workload-dependent optimum that does not exist.
 from the team's main file.
% Needs in the preamble: graphicx, booktabs, amsmath, cite, url.
% Figures are read from \EFourFigDir (default: figures/). Redefine it before the
% \input if the figures live elsewhere, e.g. \newcommand{\EFourFigDir}{e4/figures}.
% References are in e4_refs.bib.
% =============================================================================
\providecommand{\EFourFigDir}{figures}
\providecommand{\cfg}[1]{\texttt{#1}}

\section{Workload-Aware Configuration Selection}
\label{sec:e4}

The previous experiments establish that Torch-TensorRT latency depends on how
the graph is partitioned and on the shape of the request. This experiment asks
the question that follows from both: if the best compiler configuration moves
with the workload, can a cheap model learn where it moves, so that a deployment
picks its configuration from the batch size and sequence length instead of
committing to one setting for everything? Rice's algorithm selection
problem~\cite{rice1976algorithm} gives the framing, and cost-model-driven
autotuners such as AutoTVM~\cite{chen2018autotvm} show that the idea pays off
when the configuration space is large and the optimum is workload dependent.

Our answer for DistilBERT on an A100 is negative, and we report it as such. A
single tuned configuration is within 0.5\% of the best tested configuration on
every held-out shape. The two selectors we built reproduce that configuration
almost everywhere and, where they deviate, they mostly lose. What does matter
is tuning at all: the tuned configuration is 1.11$\times$ faster than default
Torch-TensorRT on the unmodified model and 1.6$\times$ to 2.3$\times$ faster
once part of the graph falls back to PyTorch. Getting to a result we trust took
three iterations, and two of them changed the experiment more than the
selector did. We describe them in the order they happened, because the
intermediate findings (a model with nothing to partition, a cluster that
cancels compile-heavy jobs, and a 20\% timing artifact) are useful to the rest
of the project on their own.

% -----------------------------------------------------------------------------
\subsection{Problem Formulation}
\label{sec:e4-problem}

A workload is a request shape $w=(b,s)$, with batch size $b$ and sequence
length $s$ in tokens. A configuration $c \in \mathcal{C}$ is one way of running
the model: PyTorch eager, or Torch-TensorRT with a particular set of compiler
and runtime settings. $L(c,w)$ is the warmed-up median latency of $c$ on $w$. A
selector is a function $\pi:\,w \mapsto c$ fitted on measurements from a set of
training shapes $W_\mathrm{train}$ and judged on shapes it has not seen. For a
test set $W$ we report the \emph{regret} of $\pi$ against the best tested
configuration,
\begin{equation}
  R(\pi) \;=\;
  \Bigl(\,\prod_{w\in W}
     \frac{L(\pi(w),w)}{\min_{c\in\mathcal{C}} L(c,w)}\Bigr)^{1/|W|} - 1 ,
  \label{eq:regret}
\end{equation}
that is, the geometric-mean slowdown relative to an oracle that tries every
configuration on every shape. A geometric mean keeps a 5\% loss at batch~1 and
a 5\% loss at batch~16 on equal footing, although the second is ten times
larger in milliseconds.

We compare six methods, following the evaluation plan of the proposal:
\begin{itemize}
  \item \textbf{PyTorch eager}, the uncompiled FP16 model;
  \item \textbf{Default Torch-TensorRT}, the compiler's default settings;
  \item \textbf{One fixed tuned configuration}, the configuration with the
        lowest geometric-mean slowdown over $W_\mathrm{train}$, used for every
        shape;
  \item \textbf{Lookup selector}, which returns the best configuration of the
        nearest training shape, with distance measured in
        $(\log_2 b,\log_2 s)$;
  \item \textbf{Regression selector}, which fits one latency model per
        configuration and returns the configuration with the lowest
        prediction (Section~\ref{sec:e4-selectors});
  \item \textbf{Best tested configuration}, the oracle in~\eqref{eq:regret}.
\end{itemize}
The fixed tuned configuration is the baseline that decides the experiment. A
selector is only worth deploying if it beats the best single configuration,
not merely the default one.

% -----------------------------------------------------------------------------
\subsection{Experimental Setup}
\label{sec:e4-setup}

All measurements come from one NVIDIA A100-SXM4-40GB node of the NYU HPC cloud
bursting cluster, inside a Singularity container with PyTorch~2.14.1
(CUDA~13.0)~\cite{paszke2019pytorch}, Torch-TensorRT~2.14.0~\cite{torchtrt} and
TensorRT~11.1.0~\cite{tensorrt}. The model is
\cfg{distilbert-base-uncased}~\cite{sanh2019distilbert} in FP16, wrapped so
that it takes token ids and an attention mask and returns the last hidden
state. Inputs are random token ids from a fixed seed with a full attention
mask, and every configuration receives the same tensors. Torch-TensorRT uses
its Dynamo front end, and each engine is built for one static shape, so a
(configuration, shape) pair is one compilation.

The training grid is $b\in\{1,4,8,16\}\times s\in\{128,256,512\}$, twelve
shapes. The proposal also lists 1024 tokens, but DistilBERT has 512 position
embeddings, so longer inputs cannot be run without changing the model. Eight
further shapes lie off the grid and are never used for fitting:
$(1,384)$, $(2,192)$, $(2,384)$, $(4,320)$, $(8,192)$, $(12,192)$, $(12,384)$
and $(16,448)$. None of their sequence lengths occurs in the grid, and four
of them also use a batch size (2 or 12) that does not.

Latency is measured with the project's shared benchmark routine: 20 warm-up
calls followed by 100 calls timed with CUDA events, from which we take the
median and the 95th percentile. Around it, our sweep harness records the
compile time, the time of the first call, the number of TensorRT and PyTorch
blocks in the compiled module, and the agreement of the output with eager
PyTorch (cosine similarity and largest absolute difference). The harness
appends one row per timing repeat to a CSV file and skips pairs that are
already present, so an interrupted job resumes where it stopped. That property
turned out to be necessary (Section~\ref{sec:e4-iter2}).

\begin{table}[t]
\caption{Chronology of cluster runs for this experiment (all on 12 Nov.\ 2026)}
\label{tab:e4-runs}
\centering
\footnotesize
\begin{tabular}{@{}llp{4.55cm}@{}}
\toprule
Job & Purpose & Outcome \\
\midrule
29499 & Environment & Cancelled: \cfg{--fakeroot} unpacked the 8.5\,GB image to
        network storage \\
29507 & Environment & Held by the scheduler; resubmitted \\
29509 & Environment & Packages installed (plain writable overlay) \\
29537 & Partition check & One engine, no fallback; LayerNorm exclusion had no
        effect (wrong operator name) \\
29550 & Partition check & Corrected; Table~\ref{tab:e4-partition} \\
29551 & Sweep 1 & GPU node never booted; cancelled after 12\,min \\
29586 & Sweep 1 & 83 of 340 runs, then cancelled by the cluster's
        low-utilization watchdog at 30\,min \\
29727--8 & Sweep 2 & Two more node boot failures (13\,min each) \\
29729 & Sweep 2 & 340 of 340 runs in 2\,h\,31\,min; all results below \\
\bottomrule
\end{tabular}
\end{table}

Table~\ref{tab:e4-runs} lists every cluster job behind this section. Only the
last one contributes numbers to the results; the others are the reason the
last one looks the way it does.

% -----------------------------------------------------------------------------
\subsection{Iteration 1: Is There Anything to Tune?}
\label{sec:e4-iter1}

Our plan assumed that the partitioning controls of Torch-TensorRT
(\cfg{min\_block\_size}, the choice between the fast and the global
partitioner, and operator exclusion) would produce visibly different graphs
for DistilBERT, and that the selector would arbitrate between them. Before
spending GPU hours on engines we checked that assumption with a dry-run
compile of every candidate at $(4,128)$, which reports the partition without
building anything.

The assumption was wrong. After lowering, DistilBERT consists of 127
operators (36 linear layers, 24 reshapes, 24 permutes, 13 additions, 13 layer
normalizations, six attention and six GELU operators, and a handful of
embedding and mask operations), and TensorRT supports all of them. The default
compile is therefore a single engine with no PyTorch fallback, and
\cfg{min\_block\_size} and the partitioner have nothing to act on: we verified
that \cfg{min\_block\_size}$\,=1$ and the global partitioner return the
identical one-engine graph, and removed both from the sweep.

That left two options: switch to a model with unsupported operators, which
would have cut this experiment off from the shared harness and from the other
three experiments, or keep the model and create the partitioning problem
deliberately. We chose the second and say so plainly, because it limits what
the results mean (Section~\ref{sec:e4-threats}). Torch-TensorRT lets the user
force named operators to run in PyTorch. We define three \emph{scenarios}:
\begin{itemize}
  \item \textbf{Unmodified model}: nothing excluded, one engine.
  \item \textbf{GELU in PyTorch}: the six GELU operators are excluded, which
        splits the encoder between its layers.
  \item \textbf{LayerNorm in PyTorch}: the thirteen layer normalizations are
        excluded, which cuts every layer twice.
\end{itemize}
Each scenario is evaluated separately, with its own default, its own tuned
configuration and its own selectors, as if it were a different model.

\begin{table}[t]
\caption{Partition structure of each configuration (dry-run compile)}
\label{tab:e4-partition}
\centering
\footnotesize
\setlength{\tabcolsep}{3.5pt}
\begin{tabular}{@{}lccl@{}}
\toprule
 & \multicolumn{2}{c}{Blocks} & \\
\cmidrule(lr){2-3}
Configuration & TRT & PyT. & Operators per TRT block \\
\midrule
\cfg{trt\_default}, \cfg{trt\_opt*} & 1 & 0 & 127 \\
\midrule
\cfg{gelu\_mbs5}, \cfg{gelu\_global} & 6 & 6 & 23, then $5\times19$ \\
\cfg{gelu\_mbs1} & 7 & 6 & 23, $5\times19$, 3 \\
\cfg{gelu\_mbs22} & 1 & 1 & 23 (104 ops in PyTorch) \\
\midrule
\cfg{ln\_mbs5} & 7 & 7 & 6, then $6\times14$ \\
\cfg{ln\_mbs1} & 13 & 13 & 6, then $6\times(14,\,4)$ \\
\cfg{ln\_mbs10} & 6 & 7 & $6\times14$ \\
\cfg{ln\_mbs20} & 0 & 1 & none (127 ops in PyTorch) \\
\cfg{ln\_global} & 6 & 7 & 18, then $5\times14$ \\
\bottomrule
\end{tabular}
\end{table}

Table~\ref{tab:e4-partition} shows the resulting partitions. The dry run also
pruned the configuration list a second time. With GELU excluded, the TensorRT
blocks hold 19 to 23 operators, so \cfg{min\_block\_size} values of 5, 10 and
20 all give the same graph; we kept 5 and replaced the other two by 22, the
smallest value that changes anything. Two mistakes surfaced in the same pass.
Our first LayerNorm scenario excluded \cfg{aten.layer\_norm}, but the operator
that exists after lowering is \cfg{aten.native\_layer\_norm}; the exclusion
silently matched nothing and the ``broken'' model was still one engine. And
our block counter reported zero PyTorch blocks for the global partitioner,
because that partitioner leaves fallback operators inline in the top-level
graph instead of wrapping them in submodules. Both were fixed before any
latency was measured. Without the dry run, a third of the sweep would have
been mislabeled.

\begin{table}[t]
\caption{Configurations in the sweep. Default settings are
\cfg{min\_block\_size}$\,=5$, fast partitioner, optimization level~3.}
\label{tab:e4-configs}
\centering
\footnotesize
\setlength{\tabcolsep}{3.5pt}
\begin{tabular}{@{}lll@{}}
\toprule
Scenario & Configuration & Change from scenario default \\
\midrule
all & \cfg{eager} & PyTorch, no compilation \\
\midrule
Unmodified & \cfg{trt\_default} & (default) \\
 & \cfg{trt\_opt1} & optimization level 1 \\
 & \cfg{trt\_opt5} & optimization level 5 \\
 & \cfg{trt\_cudagraphs} & CUDA Graphs at run time \\
 & \cfg{trt\_default\_rebuild} & none (second build) \\
\midrule
GELU in & \cfg{gelu\_mbs5} & (default) \\
PyTorch & \cfg{gelu\_mbs1} & \cfg{min\_block\_size} 1 \\
 & \cfg{gelu\_mbs22} & \cfg{min\_block\_size} 22 \\
 & \cfg{gelu\_global} & global partitioner \\
 & \cfg{gelu\_cudagraphs} & CUDA Graphs at run time \\
\midrule
LayerNorm & \cfg{ln\_mbs5} & (default) \\
in PyTorch & \cfg{ln\_mbs1} & \cfg{min\_block\_size} 1 \\
 & \cfg{ln\_mbs10}, \cfg{ln\_mbs20} & \cfg{min\_block\_size} 10, 20 \\
 & \cfg{ln\_global} & global partitioner \\
 & \cfg{ln\_cudagraphs} & CUDA Graphs at run time \\
\bottomrule
\end{tabular}
\end{table}

The final space has 17 configurations (Table~\ref{tab:e4-configs}). Because
partitioning offers nothing on the unmodified model, that scenario varies the
builder's optimization level instead. All three scenarios include whole-graph
CUDA Graphs~\cite{gray2019cudagraphs}, a runtime option that records the
module's stream of GPU kernel launches once and replays it as a single
submission. One more configuration, \cfg{trt\_default\_rebuild}, is the
default built a second time. It belongs to no scenario and exists to measure
how much two builds of the same configuration differ.

% -----------------------------------------------------------------------------
\subsection{Iteration 2: A First Sweep That Could Not Be Used}
\label{sec:e4-iter2}

The first sweep timed each compiled module three times with the shared
protocol. It was cancelled after 30 minutes with 83 of 340 runs done, and when
we looked at those 83 runs they were not usable either. The two problems are
unrelated.

\subsubsection{Compile-heavy jobs are cancelled}
The accounting record showed the job cancelled by the administrator account
with no message. The cause is a watchdog on the cluster that cancels a GPU job
once it is about 30 minutes old if its mean GPU utilization is 5\% or less. A
configuration sweep is exactly such a job: more than 85\% of its wall time is
TensorRT building engines, which is mostly CPU work, and timing a sub-millisecond engine
for 100 iterations keeps the GPU busy for a tenth of a second. We did not add
artificial load to stay above the threshold. The harness was already
resumable, so the second sweep was submitted as a chain of jobs, each starting
when the previous one ends and continuing from the same file.

\subsubsection{A 20\% timing artifact}
\label{sec:e4-warmup}

\begin{figure}[t]
\centering
\includegraphics[width=\columnwidth]{\EFourFigDir/fig_warmup.pdf}
\caption{Each timing repeat relative to the fastest repeat of the same
compiled module, for the five shapes both sweeps cover. (a)~With the shared
20-iteration warm-up, GPU-bound configurations start about 20\% slow and drop
to their final level one or two repeats later. (b)~With two seconds of
sustained load before timing, all repeats agree.}
\label{fig:e4-warmup}
\end{figure}

For 30 of the 79 successful runs, the three repeats of the \emph{same engine}
disagreed by more than 8\%. The disagreement has a clear shape
(Fig.~\ref{fig:e4-warmup}a): the first one or two repeats are 20\% to 25\%
slower, then the latency drops and stays down. It never goes the other way. It
affects only configurations whose latency is set by the GPU, namely the single
engine and everything run under CUDA Graphs; eager PyTorch and the
configurations with PyTorch fallback blocks, whose latency is set by the CPU
launching kernels, are stable to within 4\%.

Our reading is that the GPU needs a few hundred milliseconds of sustained
load to reach full speed after sitting idle through a compile, and that the
20-iteration warm-up is over long before that when an iteration takes
0.6\,ms. We did not confirm the mechanism (clock or power state) directly, so
we treat it as an observation about this node. The effect on a selector study
is severe either way. The differences we are trying to resolve on the
unmodified model are of the order of 10\%, and with three repeats the median
lands on the slow or the fast level essentially at random. Under these
conditions the selector would have been fitted to noise, and it would have
looked as if the best configuration changed from shape to shape. That
steady-state measurement is easy to get wrong without doing anything obviously
wrong is well documented~\cite{georges2007statistically,mytkowicz2009wrong};
this is one more instance.

% -----------------------------------------------------------------------------
\subsection{Iteration 3: The Corrected Sweep}
\label{sec:e4-iter3}

The second sweep differs from the first in three ways. Before timing, the
harness now keeps the GPU busy with the workload itself for two seconds; the
shared benchmark routine then runs unchanged, including its own 20 warm-up
calls. Each module is timed five times instead of three, and the latency we
use is the median of the five medians. And every row records the node and job
it came from, in case a chained sweep ends up spread over several machines.
The first sweep's data was discarded.

After two more failed node boots, the third job of the chain ran all 340 pairs
in one piece, in 2\,h\,31\,min on a single node. The warm-up had an unplanned
side effect: it lifted the job's GPU utilization above the watchdog threshold,
so the chain was never needed.

The data is clean. Of the 340 pairs, 320 succeeded. The 20 failures are all
\cfg{ln\_global}, which fails to build on every shape with an internal error
in Torch-TensorRT 2.14 (a missing \cfg{\_frozen\_param0} attribute on the
partitioned graph); it is excluded from selection. Across the five repeats the
median spread is 0.7\%, the 90th percentile 2.4\% and the maximum 5.8\%
(Fig.~\ref{fig:e4-warmup}b). The default configuration and its rebuild differ
by 0.3\% in the median. We therefore treat differences below about 1\% as
ties. Every compiled module agrees with eager PyTorch to a cosine similarity
of at least 0.99999, with a largest absolute difference of 0.02 in FP16.

% -----------------------------------------------------------------------------
\subsection{Selectors}
\label{sec:e4-selectors}

Both selectors are deliberately small. With twelve training shapes there is no
room for anything else.

The \emph{lookup selector} stores, for each training shape, the configuration
that was fastest there, and answers a query with the entry of the nearest
training shape in $(\log_2 b, \log_2 s)$. It has no parameters and cannot
extrapolate.

The \emph{regression selector} fits, for each configuration $c$, the latency
model
\begin{equation}
  \hat{L}_c(b,s) \;=\; \theta_{c,0} + \theta_{c,1}\,b + \theta_{c,2}\,bs
                      + \theta_{c,3}\,bs^2 ,
  \qquad \theta_c \ge 0 ,
  \label{eq:model}
\end{equation}
whose terms stand for a fixed per-call overhead, a per-sequence cost, a
per-token cost and the quadratic cost of self-attention. The coefficients are
obtained by non-negative least squares~\cite{lawson1995solving} on the
\emph{relative} error, $\sum_{w}(\hat{L}_c(w)/L(c,w)-1)^2$, so that small and
large shapes carry equal weight. Non-negativity keeps the model physically
meaningful and prevents it from predicting negative latencies outside the
grid. The selector returns $\arg\min_c \hat{L}_c(b,s)$.

We evaluate in two ways. \emph{Held-out}: fit on the twelve grid shapes, test
on the eight off-grid shapes. \emph{Leave-one-out}: for each grid shape, fit
on the other eleven and test on the one left out. To relate quality to tuning
effort, we also refit every method on subsets of $k=1,\dots,12$ grid shapes
(all subsets where there are at most 40, otherwise 40 random ones) and test on
the held-out shapes.

% -----------------------------------------------------------------------------
\subsection{Results}
\label{sec:e4-results}

\subsubsection{The latency landscape}

\begin{figure*}[t]
\centering
\includegraphics[width=\textwidth]{\EFourFigDir/fig_latency.pdf}
\caption{Median latency against batch size at 512 tokens. Eager PyTorch is
flat until the largest batch, which shows that it is limited by per-call
overhead and not by GPU work. A partitioned model inherits part of that
overhead (circles, squares); CUDA Graphs removes it (triangles). All curves meet at
batch 16, where GPU work dominates.}
\label{fig:e4-latency}
\end{figure*}

Fig.~\ref{fig:e4-latency} shows the structure that everything else follows
from. Eager PyTorch takes between 3.2 and 4.1\,ms on every shape below
$(16,448)$, almost independent of size: at these sizes an FP16 DistilBERT
forward pass on an A100 is bounded by the time it takes Python and the
framework to issue a long sequence of small kernels, not by the kernels
themselves. The single
TensorRT engine removes that overhead and is 2.68$\times$ faster than eager in
the geometric mean over all 20 shapes, from 6.6$\times$ at $(1,128)$ down to
1.07$\times$ at $(16,512)$, where the GPU is finally the bottleneck.

A partitioned model sits in between, and where it sits depends on how many
times execution returns to PyTorch. With GELU excluded (six fallback blocks)
the default configuration is 1.63$\times$ slower than the single engine; with
LayerNorm excluded (seven) it is 2.49$\times$ slower and only 7\% faster than
not compiling at all. Table~\ref{tab:e4-knobs} summarizes what each control is
worth.

\begin{table}[t]
\caption{Effect of each control on latency (geometric mean over all 20 shapes,
range over shapes in parentheses)}
\label{tab:e4-knobs}
\centering
\footnotesize
\begin{tabular}{@{}p{4.35cm}l@{}}
\toprule
Change & Latency \\
\midrule
\multicolumn{2}{@{}l}{\emph{Unmodified model, relative to default}} \\
Optimization level 5 & 1.00$\times$ (0.97--1.02) \\
Optimization level 1 & 1.06$\times$ slower (1.03--1.08) \\
CUDA Graphs & 1.11$\times$ faster (1.02--1.20) \\
\midrule
\multicolumn{2}{@{}l}{\emph{GELU in PyTorch, relative to its default}} \\
\cfg{min\_block\_size} 1 (7 engines) & 1.01$\times$ faster (0.99--1.04) \\
Global partitioner & 1.02$\times$ slower (0.98--1.05) \\
\cfg{min\_block\_size} 22 (1 engine) & 1.74$\times$ slower (1.07--2.25) \\
CUDA Graphs & 1.63$\times$ faster (1.06--3.27) \\
\midrule
\multicolumn{2}{@{}l}{\emph{LayerNorm in PyTorch, relative to its default}} \\
\cfg{min\_block\_size} 1 (13 engines) & 1.05$\times$ faster (0.99--1.11) \\
\cfg{min\_block\_size} 10 (6 engines) & 1.03$\times$ slower (1.00--1.05) \\
\cfg{min\_block\_size} 20 (no engine) & 1.28$\times$ slower (1.04--1.47) \\
CUDA Graphs & 2.26$\times$ faster (1.08--5.02) \\
\bottomrule
\end{tabular}
\end{table}

Two things stand out. First, the partitioning controls barely matter as long
as the engines survive: \cfg{min\_block\_size} and the partitioner move
latency by 5\% or less, and their only large effect is the destructive one,
when a high threshold discards engines and pushes the model back towards
eager. Second, one runtime option dwarfs all of them. CUDA Graphs makes the
GELU-partitioned model 1.63$\times$ faster and the LayerNorm-partitioned one
2.26$\times$ faster, to within 11\% and 23\% of the single engine under the
same option. This is consistent with the picture above: what a fallback block
costs is launch overhead, and replaying a recorded launch sequence removes it
no matter how many blocks there are.

\begin{figure*}[t]
\centering
\includegraphics[width=\textwidth]{\EFourFigDir/fig_gain.pdf}
\caption{Speedup of the best tested configuration over the scenario's default
on the training grid. Tuning is worth most at small shapes and when the graph
is partitioned; it fades as the GPU becomes the bottleneck.}
\label{fig:e4-gain}
\end{figure*}

Fig.~\ref{fig:e4-gain} shows the same thing from the point of view of tuning.
The \emph{size} of the gain is strongly workload dependent, from 5.0$\times$
at $(1,128)$ with LayerNorm excluded down to a few percent at $(16,512)$. The
\emph{identity} of the best configuration is not, and only the latter is
something a selector can exploit.

\subsubsection{Selection on held-out shapes}

\begin{table}[t]
\caption{Held-out shapes: regret $R$ from~\eqref{eq:regret} and speedup over
default Torch-TensorRT. GELU and LayerNorm denote the scenarios in which that
operator runs in PyTorch.}
\label{tab:e4-heldout}
\centering
\footnotesize
\setlength{\tabcolsep}{2.6pt}
\begin{tabular}{@{}lrrrrrr@{}}
\toprule
 & \multicolumn{2}{c}{Unmodified} & \multicolumn{2}{c}{GELU}
 & \multicolumn{2}{c}{LayerNorm} \\
\cmidrule(lr){2-3}\cmidrule(lr){4-5}\cmidrule(l){6-7}
Method & $R$\,(\%) & speedup & $R$\,(\%) & speedup & $R$\,(\%) & speedup \\
\midrule
PyTorch eager        & 209.8 & 0.36$\times$ & 180.9 & 0.57$\times$ & 153.8 & 0.92$\times$ \\
Default TRT          &  11.0 & 1.00$\times$ &  61.3 & 1.00$\times$ & 132.7 & 1.00$\times$ \\
Fixed tuned          &   0.0 & 1.11$\times$ &   0.0 & 1.61$\times$ &   0.5 & 2.32$\times$ \\
Lookup sel.          &   0.0 & 1.11$\times$ &   3.2 & 1.56$\times$ &   1.4 & 2.29$\times$ \\
Regression sel.      &   1.5 & 1.09$\times$ &   0.5 & 1.60$\times$ &   0.0 & 2.33$\times$ \\
Best tested          &   0.0 & 1.11$\times$ &   0.0 & 1.61$\times$ &   0.0 & 2.33$\times$ \\
\bottomrule
\end{tabular}
\end{table}

\begin{figure}[t]
\centering
\includegraphics[width=\columnwidth]{\EFourFigDir/fig_regret.pdf}
\caption{Regret on the eight held-out shapes for the three methods that use
tuning data. The dashed line is the timing noise (0.7\%). No selector beats
the fixed tuned configuration by more than the noise.}
\label{fig:e4-regret}
\end{figure}

Table~\ref{tab:e4-heldout} and Fig.~\ref{fig:e4-regret} give the main result.
In all three scenarios the fixed tuned configuration is the CUDA Graphs
variant, and it is the best tested configuration on 23 of the 24
scenario--shape combinations. Its regret is 0.0\%, 0.0\% and 0.5\%. The
selectors range from 0.0\% to 3.2\%; relative to the fixed tuned configuration
they achieve between 0.97$\times$ and 1.01$\times$. The single case in which a
selector is ahead (regression, LayerNorm scenario, by 0.5\%) is below the
noise floor. Leave-one-out evaluation on the grid agrees: the fixed tuned
configuration has a regret of 0.0\%, 0.1\% and 0.9\%, the lookup selector
makes identical choices, and the regression selector scores 0.8\%, 1.0\% and
0.0\%.

Against default Torch-TensorRT, on the other hand, the gain is real and
sizeable: 1.11$\times$ on the unmodified model, 1.61$\times$ with GELU
excluded and 2.32$\times$ with LayerNorm excluded.

\subsubsection{Why there is nothing to select}

\begin{figure}[t]
\centering
\includegraphics[width=\columnwidth]{\EFourFigDir/fig_choice.pdf}
\caption{Fastest configuration at each of the 20 shapes in the two
partitioned scenarios. The only shape dependence is a switch to eager PyTorch
at the largest shapes. Circles mark held-out shapes where the lookup selector
chose differently from the best configuration, with the resulting loss.}
\label{fig:e4-choice}
\end{figure}

On the unmodified model, \cfg{trt\_cudagraphs} is the fastest configuration on
all 20 shapes. In the partitioned scenarios the CUDA Graphs variant is fastest
everywhere except at the top-right corner of the shape space
(Fig.~\ref{fig:e4-choice}): at $(16,512)$ in both scenarios and at $(16,448)$
with LayerNorm excluded, eager PyTorch is faster by 1\% to 11\%, because at
that size the per-call overhead no longer matters and the compiled module
pays for crossing between engines and PyTorch.

That one crossover is the entire workload dependence in our data, and the
selectors handle it poorly. Across the three scenarios the two selectors chose
eager for a held-out shape seven times, and it was the right call twice. The
lookup selector assigns $(12,384)$ to its nearest grid neighbor $(16,512)$ and
therefore to eager, at a cost of 23\% on that shape with GELU excluded and
12\% with LayerNorm excluded. The regression selector does better on the
partitioned scenarios but sends $(16,448)$ to eager on the unmodified model,
where the engine is 13\% faster. The underlying problem is that the boundary
is defined by a single training point at the edge of the grid, and that the
latency model is least accurate for exactly the configuration on the other
side of it: its mean absolute error on held-out shapes is 3\% to 4\% for the
GPU-bound configurations but 13\% for eager, whose flat-then-rising curve the
four terms of~\eqref{eq:model} do not fit well.

\subsubsection{Tuning cost}

\begin{figure*}[t]
\centering
\includegraphics[width=\textwidth]{\EFourFigDir/fig_budget.pdf}
\caption{Held-out regret when only $k$ of the twelve grid shapes are measured
(mean over subsets). Two shapes are enough to identify the fixed tuned
configuration. With one shape the outcome depends on which: if it is
$(16,512)$, every method selects eager PyTorch.}
\label{fig:e4-budget}
\end{figure*}

\begin{table}[t]
\caption{Tuning cost and break-even against default Torch-TensorRT, per
scenario}
\label{tab:e4-cost}
\centering
\footnotesize
\setlength{\tabcolsep}{4pt}
\begin{tabular}{@{}lrrr@{}}
\toprule
 & Unmodified & GELU & LayerNorm \\
\midrule
Saving per request (ms)        & 0.12 & 0.70 & 1.76 \\
\midrule
Full grid, 12 shapes (GPU-min) & 16.6 & 30.1 & 33.3 \\
\quad break-even (requests)    & $8.3\times10^6$ & $2.6\times10^6$ & $1.1\times10^6$ \\
\midrule
Two shapes (GPU-min)           & 2.7 & 5.1 & 5.5 \\
\quad break-even (requests)    & $1.4\times10^6$ & $4.3\times10^5$ & $1.9\times10^5$ \\
\bottomrule
\end{tabular}
\end{table}

The cost of a measurement is the engine build, the first call, the warm-up
and the timing. Measuring every configuration of a scenario on all twelve
grid shapes costs 17 to 33 GPU-minutes (Table~\ref{tab:e4-cost}), of which
engine builds are the bulk: a single-engine build takes 8 to 14\,s at the
default optimization level and 32\,s at level~5, a partitioned one about
30\,s, and the 13-engine \cfg{ln\_mbs1} about 57\,s.
Fig.~\ref{fig:e4-budget} shows that most of this effort is unnecessary. Two
measured shapes identify the fixed tuned configuration in every scenario, and
the ten further shapes buy nothing for it; the selectors need more data and
still end up level or behind. Dividing the cost by the mean saving per request
on the held-out shapes gives the number of requests after which tuning has
paid for itself: about 1.4 million on the unmodified model, where the saving
is 0.12\,ms, and 0.2 to 0.4 million in the partitioned scenarios.

% -----------------------------------------------------------------------------
\subsection{Discussion}
\label{sec:e4-discussion}

The proposal's hypothesis for this experiment was that a lookup table or a
regression model would recover most of the gap between a fixed configuration
and the best one. For this model the hypothesis fails for an uninteresting
reason: the gap is not there. We think the reason generalizes somewhat beyond
our setup, and it is worth stating.

A selector needs configurations whose latency curves \emph{cross}. In our
space, the configurations differ in how much per-call overhead they carry, and
overhead is a constant offset: a configuration with less of it is faster at
every shape, by an amount that shrinks as the shape grows but does not change
sign. CUDA Graphs removes the most overhead and therefore wins everywhere it
is applicable. The only crossing we observe involves eager PyTorch at the
largest shapes, and it is shallow (at most 11\%) and located at the edge of
the measured region. Per-workload selection becomes interesting when
configurations trade a fixed cost against a marginal one, for instance
precision modes with different kernel speeds, engines built with dynamic
versus static shape profiles, or models where a fallback block does real
computation instead of just interrupting the engine. None of these is in our
space.

For the project as a whole, three findings from this experiment matter more
than the selector:
\begin{enumerate}
  \item DistilBERT converts to a single TensorRT engine. Any partitioning
        effect studied on it has to be induced.
  \item On this model the cost of a partition boundary is launch overhead,
        and CUDA Graphs recovers most of it. Whether a graph is ``well
        partitioned'' matters far less with that option enabled.
  \item The shared benchmark's 20-iteration warm-up is too short for
        sub-millisecond engines on this hardware and overstates their latency
        by about 20\% in some repeats. Experiments that compare a single
        engine against PyTorch with it will understate the speedup.
\end{enumerate}

% -----------------------------------------------------------------------------
\subsection{Threats to Validity}
\label{sec:e4-threats}

\emph{The partitioned scenarios are artificial.} GELU and LayerNorm convert
without difficulty, and nobody would exclude them in practice. The scenarios
reproduce the structure of a fallback (execution leaves the engine and comes
back) but not necessarily its cost profile, since a genuinely unsupported
operator may be expensive in itself.

\emph{One model, one GPU, one precision.} Everything here is DistilBERT in
FP16 on one A100 node, with sequences of at most 512 tokens. All 340 pairs
were measured in a single job on a single machine, which removes
machine-to-machine variation from the comparison but also means we have not
observed it.

\emph{Static shapes.} Every engine is built for exactly one shape. A
deployment that serves many shapes from one engine with a dynamic profile
faces a different trade-off, which this experiment does not measure.

\emph{Drift in the CPU-bound baselines.} The held-out shapes were measured
in the last hour of the job, and the CPU-bound configurations look slower
there than on the grid. Relative to the grid-fitted model~\eqref{eq:model},
eager PyTorch is 15\% slower on the held-out shapes on average, the
LayerNorm-partitioned configurations 8\% to 11\%, and the GPU-bound
configurations 1\%. Some of this is model error, but neighboring shapes point
the same way (eager takes 3.25\,ms at $(8,256)$ on the grid and 4.06\,ms at
$(8,192)$ held out), so host-side performance probably degraded during the
job. Two consequences follow. The held-out speedups over eager and over the
partitioned defaults are somewhat inflated; the leave-one-out figures on the
grid (2.22$\times$ over the default for LayerNorm, against 2.32$\times$ held
out) are the more conservative ones. And the position of the eager crossover
on the held-out shapes is uncertain: if eager were in fact 15\% faster at
$(16,448)$, it would be the best choice there in all three scenarios, the
fixed tuned configuration would lose 2\% to 20\% on that one shape, and a
selector choosing eager there would be right. This would move the
held-out regret of the fixed tuned configuration to at most 2.3\%. It would not
change the picture at the other seven shapes or on the grid.

\emph{Build times share a cache.} TensorRT reuses layer timings within a
process, so the first build of the job took 40\,s and later builds of the same
graph 6 to 13\,s. The tuning costs in Table~\ref{tab:e4-cost} are therefore
optimistic for a cold start.

\emph{Memory.} The recorded peak memory counts PyTorch's allocator only and
excludes memory held by TensorRT engines, so we do not report a memory
comparison.

% -----------------------------------------------------------------------------
\subsection{Summary}
\label{sec:e4-summary}

We built a resumable sweep harness, two workload-aware selectors and an
evaluation on held-out shapes, and measured 17 configurations on 20 shapes.
For DistilBERT on an A100, tuning the configuration once is worth
1.11$\times$ to 2.32$\times$ over default Torch-TensorRT, two measured shapes
are enough to do it, and selecting the configuration per workload adds nothing
beyond that. The result rests on a measurement protocol that had to be
corrected first; with the original warm-up the same experiment would have
reported a workload-dependent optimum that does not exist.
